\providecommand{\toplevelprefix}{../..}  %
\documentclass[../../book-main_zh.tex]{subfiles}

\begin{document}

\chapter{学习线性和独立结构}
\label{ch:classic-models}

\begin{quote}
    \hfill  ``{\em 做数学的艺术在于找到那个包含所有一般性萌芽的特例}。''

    $~$\hfill -- 大卫·希尔伯特 (David Hilbert)
\end{quote}
\vspace{5mm}

\textit{真实数据具有低维结构。}为了理解为什么这是真的，让我们考虑一个不起眼的例子：当卫星信号中断时电视机上的雪花噪点。在每一帧（大约每
\(\frac{1}{30}\) 秒）中，尺寸为 \(H \times W\) 的屏幕上的 RGB 噪点，大致是从 \([0, 1]^{3
\times H \times W}\)
上的均匀分布中独立采样的。理论上，这些噪点在任何给定帧中\textit{都有可能}组合成一幅自然图像，但即使你盯着电视屏幕看上一千年，这种情况也不会发生。这种差异的原因在于，\(H
\times W\) 的自然图像集合在超立方体 \([0, 1]^{3 \times H \times W}\)
中所占的比例微乎其微。特别是，与环境空间的维度相比，它是极度低维的。类似的现象也存在于所有其他类型的自然数据中，如文本、音频和视频。因此，当我们设计系统和方法来处理自然数据并学习其结构或分布时，这是我们需要考虑的自然数据的一个核心属性。
%

因此，我们的核心任务是学习一个在高维空间中具有低内在维度的分布。在本节的其余部分，我们将讨论几种执行此任务的\textit{经典}方法，这些方法针对几种多少有些{\em
理想化}的分布模型，即几何上线性或统计上独立的模型。虽然这些模型和方法本身就非常重要且有用，但我们在这里讨论它们，是因为它们能够激发、启发并作为更现代的方法（涉及深度表征学习的更一般分布的方法）的前身或类比。

我们的主要方法（以及一般的问题表述）可以总结为：
\begin{tcolorbox}\centering
    \textbf{问题：}\textit{给定来自数据分布的某个真实样本的一个或多个（带噪或不完整的）观测，对该样本进行估计。}
\end{tcolorbox}
这种方法构成了数据处理中几种经典方法的基础，我们将在本章中讨论这些方法。
\begin{itemize}
    \item \Cref{sec:lowrank} --- 主成分分析 (PCA)：给定支撑在{\em
        一个低维子空间}上的分布的嘈杂样本，获得位于该子空间上的真实样本的估计值。
    \item  \Cref{sec:ica} --- 完备字典学习与独立成分分析 (ICA)：给定支撑在{\em
        几个低维子空间的并集（\textit{而非}张成空间）}上的分布的嘈杂样本，获得真实样本的估计值。
    \item \Cref{sec:dictionary_learning} --- 稀疏编码与字典学习：给定来自{\em
        支撑在几个不相干向量（如坐标轴）组合上的分布}的嘈杂样本，获得同样具有此性质的真实样本的估计值。
\end{itemize}

如上所述，在本章中，我们做出了简化的建模假设，本质上假设数据分布（的支撑集）具有几何上（近似、分段）线性的结构和统计上独立的成分，如图
\ref{fig:subspaces} 所示。
\begin{figure}[t]
    \centering
    \includegraphics[height=0.3\linewidth]{\toplevelprefix/chapters/chapter2/figs/One-subspace.png}
    \hspace{6mm}
    \includegraphics[height=0.3\linewidth]{\toplevelprefix/chapters/chapter2/figs/subspaces.png}
    \caption{单个低维子空间上的数据（左），例如 $\mathcal{S}
        = \mathop{\mathrm{col}}(\vU)$，或低维子空间的混合（右），例如 $\mathcal{S}_j
    = \mathop{\mathrm{col}}(\vU_j)$。}
    \label{fig:subspaces}
\end{figure}
在 \Cref{ch:intro}
中，我们将此类数据模型称为“分析模型”。这些建模假设使我们能够推导出具有可证明效率保证\footnote{在数据和计算复杂度方面。}的高效算法，以大规模处理数据。然而，对于通常复杂的现实世界数据分布而言，它们是不完美的模型，因此其潜在假设仅近似成立。这意味着对这些算法进行详细分析所提供的保证在真实数据的情况下也仅近似成立。

尽管如此，本章讨论的技术本身就很有用，除此之外，它们还可以说是“包含一般性萌芽的特例”，因为它们为我们稍后要讨论的更一般分布的（深度）学习中更一般的范式提供了指导动机和直觉。正如我们将在后面的章节中很快揭示的那样，在深度学习时代，学习（低维）数据分布的现代“数据驱动”或“非参数”方法本质上采用了相同的概念和思想来进行学习。

\section{低维子空间} \label{sec:lowrank}

\subsection{主成分分析 (PCA)}\label{sub:pca}

对于低维结构，可能最简单的建模假设就是所谓的\textit{低秩}假设。设 \(D\) 为数据空间的维度，我们假设数据属于一个维度为
\(d \ll D\) 的低维子空间，可能还加上一些小的扰动。对于一些令人惊讶的复杂数据（如
\Cref{fig:faces-digits} 所示的手写数字图像和人脸数据
\cite{BasriR2003-PAMI}），这最终成为一个几乎有效的假设，而且正如我们将看到的，它非常适合进行全面的分析。

\begin{figure}
    \centering
    \includegraphics[width=0.4\linewidth]{\toplevelprefix/chapters/chapter2/figs/faces.png}
    \hspace{5mm}
    \includegraphics[width=0.395\linewidth]{\toplevelprefix/chapters/chapter2/figs/handwritten-digits.png}
    \caption{人脸和手写数字的图像。尽管它们的外观看起来变化很大，但每组数据都（近似地）张成了一个非常低维的（近似）线性子空间。}
    \label{fig:faces-digits}
\end{figure}

\paragraph{问题表述。}
为了用数学符号表示这一点，我们用一个正交矩阵 \(\vU \in \O(D, d) \subseteq \R^{D \times d}\)
来表示维度为 \(d\) 的子空间 \(\cS \subseteq \R^{D}\)，使得 \(\vU\) 的列张成
\(\cS\)。那么，如果存在一个正交矩阵 \(\vU \in \O(D, d)\)、向量 \(\{\vz_{i}\}_{i =
1}^{N} \subseteq \R^{d}\) 以及\textit{较小}的向量 \(\{\veps_{i}\}_{i =
1}^{N} \subseteq \R^{D}\)，使得
\begin{equation}\label{eq:pca_dgp}
    \vx_{i} = \vU\vz_{i} + \veps_{i}, \quad \forall i \in [N].
\end{equation}
我们就说我们的数据 \(\{\vx_{i}\}_{i = 1}^{N} \subseteq \R^{D}\) 具有（近似）低秩结构。这里的
\(\veps_{i}\)
是阻碍数据成为完美低秩的扰动；它们在模型中的存在使我们能够量化当存在偏离模型的情况时，我们的分析在多大程度上仍然具有相关性。真实的支撑子空间是
\(\cS \doteq
\mathop{\mathrm{col}}(\vU)\)，即 $\vU$
的列的张成空间。为了从这些数据中处理出我们所能得到的一切，我们需要恢复 \(\cS\)；为此，恢复
\(\vU\)（也称为\textit{主成分}）就足够了。幸运的是，这是一个计算上易于处理的任务，名为{\bf 主成分分析}，我们现在讨论如何解决它。

给定如 \eqref{eq:pca_dgp} 分布的数据 \(\{\vx_{i}\}_{i = 1}^{N}\)，我们的目标是恢复模型
\(\vU\)。一种自然的方法是找到子空间 \(\vU^{\star}\) 和潜在向量 \(\{\vz_{i}^{\star}\}_{i
= 1}^{N}\)，它们能产生最佳近似 \(\vx_{i} \approx
\vU^{\star}\vz_{i}^{\star}\)。也就是说，我们的目标是解决以下问题
\begin{equation}\label{eq:pca_sparse_recovery_problem}
    \min_{\tilde{\vU}, \{\tilde{\vz}_{i}\}_{i =
    1}^{N}}\frac{1}{N}\sum_{i = 1}^{N}\norm{\vx_{i} -
    \tilde{\vU}\tilde{\vz}_{i}}_{2}^{2},
\end{equation}
其中 $\tilde{\vU}$ 被约束为正交矩阵，如上所述。为了简洁起见，我们将在下面类似的表述中省略此约束。

\paragraph{通过去噪进行子空间编码-解码。}
为了简化这个问题，对于固定的 \(\tilde{\vU}\)，我们有（证明留作练习）：
\begin{align}
    \min_{\{\tilde{\vz}_{i}\}_{i = 1}^{N}}\frac{1}{N}\sum_{i =
    1}^{N}\norm{\vx_{i} - \tilde{\vU}\tilde{\vz}_{i}}_{2}^{2}
    &= \frac{1}{N}\sum_{i = 1}^{N}\min_{\tilde{\vz}_{i}}\norm{\vx_{i}
    - \tilde{\vU}\tilde{\vz}_{i}}_{2}^{2} \\
    &= \frac{1}{N}\sum_{i = 1}^{N}\norm{\vx_{i} -
    \tilde{\vU}\tilde{\vU}^{\top}\vx_{i}}_{2}^{2}.
\end{align}
也就是说，上述优化问题的最优解 \((\vU^{\star}, \{\vz_{i}^{\star}\}_{i = 1}^{N})\) 满足
$\vz_i^{\star} = (\vU^{\star})^\top\vx_i$。

现在，我们可以将关于 \(\tilde{\vU}^{\star}\) 和 \(\{\tilde{\vz}_{i}\}_{i =
1}^{N}\) 的原始优化问题写成仅关于 \(\tilde{\vU}\) 的优化问题，即为了获得基 \(\vU^{\star}\)
和紧凑编码 \(\{\vz_{i}^{\star}\}_{i = 1}^{N}\)，只需解决以下两个等价问题之一即可：
\begin{equation}\label{eq:pca_equals_denoising}
    \min_{\tilde{\vU}, \{\tilde{\vz}_{i}\}_{i =
    1}^{N}}\frac{1}{N}\sum_{i = 1}^{N}\norm{\vx_{i} -
    \tilde{\vU}\tilde{\vz}_{i}}_{2}^{2} =
    \min_{\tilde{\vU}}\frac{1}{N}\sum_{i = 1}^{N}\norm{\vx_{i} -
    \tilde{\vU}\tilde{\vU}^{\top}\vx_{i}}_{2}^{2}.
\end{equation}
请注意，\eqref{eq:pca_equals_denoising}
右侧的问题是一个\textit{去噪}问题：给定低秩数据的嘈杂观测值 \(\vx_{i}\)，我们的目标是找到 \(\vx_{i}\)
的\textit{无噪}副本，在模型 \eqref{eq:pca_dgp} 下即为 $\vU\vz_{i}$。也就是说，去噪后的输入为
$\hat{\vx}_i = \vU\vU^\top \vx_i$。请注意，这是子空间上距离 $\vx_i$ 最近的点，如
\Cref{fig:pca-geometry} 所示。在这里，通过解决寻找最佳子空间的等价问题（由学习到的基
\(\vU^{\star}\) 参数化），我们学习到了\textit{去噪器}的一个近似，即投影矩阵
\(\vU^{\star}(\vU^{\star})^{\top} \approx \vU\vU^{\top}\)，它将嘈杂的数据点
$\vx_i$ 投影到子空间 \(\cS\) 上。 %

\begin{figure}
    \centering
    \begin{tikzpicture}
        \node (zero) at (0, 0) {};
        \node (u1) at (5, 1) {\(\vu_{1}\)};
        \draw[blue, ->, thick] (zero) -- (u1);
        \node[circle, inner sep=1.25pt, fill=red, draw=black] (x) at
        (3, 1.5) {};
        \node[circle, inner sep=1.25pt, fill=green, draw=black]
        (x_hat) at (3.2, 0.64)  {};
        \node at ($(x) + (0.5, 0.0)$) {\(\vx\)};
        \node at ($(x_hat) + (1, -0.2)$) {\(\hat{\vx} =
        \vu_{1}\vu_{1}^{\top}\vx\)};
        \draw[brown, ->, thick] (x_hat) -- (x);
        \node at (2.75, 1) {\color{brown}\(\varepsilon\)};
    \end{tikzpicture}
    \caption{\small \textbf{PCA 的几何意义。} 数据点 $\vx$（红色）被投影到由单位基向量
        $\vu_1$（蓝色箭头）张成的一维学习子空间上。投影 $\vU\vU^\top\vx =
        \vu_{1}\vu_{1}^{\top}\vx$（绿色）是利用 $\vu_{1}$ 给出的低维结构对 $\vx$
    进行去噪后的版本，而 $\boldsymbol{\varepsilon}$（棕色箭头）表示投影残差或噪声。}
    \label{fig:pca-geometry}
\end{figure}

将上述过程结合起来，我们本质上获得了一个简单的编码-解码方案，该方案将 $\R^D$ 中的数据点 $\vx$ 映射到低维（潜在）空间
$\R^d$，然后再映射回 $\R^D$：
\begin{equation}
    \x \xrightarrow{\hspace{2mm} \mathcal{E} = (\vU^\star)^\top
    \hspace{2mm}}  \z
    \xrightarrow{\hspace{2mm} \mathcal{D} = \vU^\star \hspace{2mm}}   \hat{\x}.
    \label{eqn:autoencoding-PCA}
\end{equation}
在这里，$\vz \in \R^d$ 可以被视为数据点 $\vx \in \R^D$ 的低维紧凑编码（或潜在表征），而学习到的子空间基
$\vU^\star$ 可以被视为相关的码本，其列是（学习到的）最优码字。该过程通过将 $\vx$ 投影到由 $\vU^\star$
张成的子空间上，实现了对 $\vx$ 的去噪功能。

\paragraph{计算子空间基。}
现在，我们继续计算。设 \(\vX = \mat{\vx_{1},
\dots, \vx_{N}} \in \R^{D \times N}\) 为以观测值 \(\vx_{i}\) 为列的矩阵。我们有（证明留作练习）
\begin{align}
    \argmin_{\tilde{\vU}}\frac{1}{N}\sum_{i = 1}^{N}\norm{\vx_{i} -
    \tilde{\vU}\tilde{\vU}^{\top}\vx_{i}}_{2}^{2}
    &= \argmax_{\tilde{\vU}}\frac{1}{N}\sum_{i =
    1}^{N}\norm{\tilde{\vU}^{\top}\vx_{i}}_{F}^{2} \\
    &=
    \argmax_{\tilde{\vU}}\tr\rc{\tilde{\vU}^{\top}\bp{\frac{\vX\vX^{\top}}{N}}\tilde{\vU}}.
\end{align}
因此，为了计算主成分，我们寻找使项
\(\tr(\tilde{\vU}^{\top}(\vX\vX^{\top}/N)\tilde{\vU})\) 最大化的正交矩阵
\(\tilde{\vU}\)。我们可以通过数学归纳法证明，该矩阵 \(\vU^{\star}\) 的列是
\(\vX\vX^{\top}/N\) 的\textit{前 \(d\) 个单位特征向量}。我们将完整的证明留给读者在
\Cref{exercise:principal-components-derivation}
中完成，但我们在这里处理归纳法的基础情况。假设 \(d = 1\)。那么我们只需要恢复一个单位向量 \(\vu\)，因此上述问题简化为
\begin{equation}
    \max_{\tilde{\vu} \colon \norm{\tilde{\vu}}_{2} = 1}
    \tilde{\vu}^{\top}(\vX\vX^{\top}/N)\tilde{\vu}.
\end{equation}
这就是所谓的 \(\vX\vX^{\top}/N\) 的\textit{瑞利商 (Rayleigh
quotient)}。通过调用谱定理，我们将 \(\vX\vX^{\top}/N\) 对角化为 \(\vX\vX^{\top}/N =
\vV\vLambda\vV^{\top}\)，其中 \(\vV\) 是正交矩阵，\(\vLambda\) 是具有非负元素的对角矩阵。因此
\begin{equation}
    \tilde{\vu}^{\top}(\vX\vX^{\top}/N)\tilde{\vu} =
    \tilde{\vu}^{\top}\vV\vLambda\vV^{\top}\vu =
    (\vV^{\top}\tilde{\vu})^{\top}\vLambda(\vV^{\top}\tilde{\vu}).
\end{equation}
由于 \(\vV\) 是可逆的正交变换，\(\vV^{\top}\tilde{\vu}\) 是一个单位向量，并且对
\(\tilde{\vu}\) 进行优化等价于对 \(\tilde{\vw} \doteq
\vV^{\top}\tilde{\vu}\) 进行优化。因此，我们可以写成
\begin{equation}
    \tilde{\vu}^{\top}(\vX\vX^{\top}/N)\tilde{\vu} =
    \tilde{\vw}^{\top}\vLambda\tilde{\vw},
\end{equation}
在单位向量中，其最优解 \(\vw^{\star}\) 是独热 (one-hot) 向量，其唯一非零（因此为单位值）的元素位于对应于
\(\vX\vX^{\top}/N\) 最大特征值的索引之一。这意味着原始问题的最优解 \(\tilde{\vu} =
\vV\tilde{\vw}\) 对应于 \(\vX\vX^{\top}/N\) 的一个单位特征向量（即 \(\vV\)
的列），该特征向量对应于最大特征值。将其适当地推广到 \(d > 1\) 的情况，并总结前面的所有讨论，我们得到以下非正式定理。
\begin{theorem}\label{thm:pca}
    假设我们的数据集 \(\{\vx_{i}\}_{i = 1}^{N} \subseteq
    \R^{D}\) 可以写成
    \begin{equation}
        \vx_{i} = \vU\vz_{i} + \veps_{i}, \qquad \forall i \in [N],
    \end{equation}
    其中 \(\vU \in \O(D, d)\) 捕获了\textit{低秩结构}，\(\{\vz_{i}\}_{i =
    1}^{N} \subseteq \R^{d}\) 是数据的\textit{紧凑编码}，而 \(\{\veps_{i}\}_{i
    = 1}^{N} \subseteq \R^{D}\)
    是表示我们的数据偏离低秩模型的小向量。那么我们数据集的\textit{主成分} \(\vU^{\star} \in \O(D,
    d)\) 由 \(\vX\vX^{\top}/N\) 的前 \(d\) 个特征向量给出，其中 \(\vX = [\vx_{1},
        \dots, \vx_{N}] \in \R^{D \times
    N}\)，并且近似对应于最优线性去噪器：\(\vU^{\star}(\vU^{\star})^{\top} \approx
    \vU\vU^{\top}\)。
\end{theorem}
我们在这里不给出明确的近似率，因为它们可能会变得相当技术化。在所有 \(i\) 都有 \(\veps_{i} = \vzero\)
的特殊情况下，学习到的 \(\vU^{\star}\) 张成了样本 \(\{\vx_{i}\}_{i
= 1}^{N}\) 的支撑集。如果此外 \(\vz_{i}\) 足够多样化（例如，张成整个
\(\R^{d}\)），那么我们将获得完美恢复：\(\vU^{\star}(\vU^{\star})^{\top} = \vU\vU^{\top}\)。

\begin{remark}
    在某些数据分析任务中，数据矩阵 \(\vX\)
    的格式使得每个数据点是一\textit{行}，而不是这里所呈现的一\textit{列}。在这种情况下，主成分是
    \(\vX^{\top}\vX/N\) 的前 \(d\) 个特征向量。
\end{remark}

\begin{remark}[通过去噪特征值进行基选择]
    在许多情况下，要么我们的数据不会真正按照子空间加噪声的模型分布，要么我们不知道子空间真实的潜在维度
    \(d\)。在这种情况下，我们必须选择 \(d\)；这个问题被称为\textit{模型选择}。在 PCA
    的受限情况下，执行模型选择的一种方法是计算 \(\vX\vX^{\top}/N\)
    并寻找相邻特征值急剧下降的情况；这是一个指标，表明较大特征值的索引是“真实维度 \(d\)”，而
    \(\vX\vX^{\top}/N\) 的其余特征值是由噪声或扰动 \(\veps_{i}\)
    贡献的。模型选择是一个困难的问题，如今在深度学习时代，它被称为“超参数优化”，通常通过暴力搜索或贝叶斯优化来完成，通常使用保留（“验证”）数据集上的损失作为目标。
    %
\end{remark}

\begin{remark}[去噪样本]
    \eqref{eq:pca_equals_denoising} 右侧的表达式，即
    \begin{equation}\label{eq:orthogonal_denoising}
        \min_{\tilde{\vU}}\frac{1}{N}\sum_{i = 1}^{N}\norm{\vx_{i} -
        \tilde{\vU}\tilde{\vU}^{\top}\vx_{i}}_{2}^{2},
    \end{equation}
    就是所谓的\textit{去噪问题}，之所以这样命名，是因为它是一个优化问题，其解\textit{去除了样本中的噪声，使其拟合在子空间上}。去噪——学习一个映射，从嘈杂样本中去除噪声，使其拟合数据结构（如在
    \eqref{eq:pca_dgp}
    中，但可能更复杂）——是学习分布的一种常用方法，将在后文和整部书稿中讨论。请注意，我们已经讨论过这个概念，但由于它在后面章节中的核心重要性，值得重复强调。
\end{remark}

\begin{remark}[神经网络解释]
    如果我们进行 PCA，我们近似恢复了由参数 \(\vU^{\star}\) 编码的分布的支撑集。学习到的去噪映射随后采用
    \(\vU^{\star}(\vU^{\star})^{\top}\)
    的形式。除了作为去噪器之外，这还可以被视为一个\textit{简单的两层权重绑定线性神经网络}：第一层乘以
    \((\vU^{\star})^{\top}\)，第二层乘以 \(\vU^{\star}\)，即
    \begin{equation}
        \operatorname{denoise}(\vx) = \underbrace{\vU^{\star} \circ
            \underbrace{\id \circ
        \underbrace{(\vU^{\star})^{\top}\vx}_{\text{第一“层”}}}_{\text{第一“层”的激活后}}}_{\text{“神经网络”的输出}}
    \end{equation}
    将其与标准的双层神经网络进行对比，我们看到了结构上的相似性：
    \begin{equation}
        \operatorname{NN}(\vx) = \underbrace{\vW^{\star} \circ
            \underbrace{\mathrm{ReLU} \circ
        \underbrace{(\vU^{\star})^{\top}\vx}_{\text{第一层}}}_{\text{第一层的激活后}}}_{\text{神经网络的输出}}
    \end{equation}
    特别是，PCA
    可以被解释为\textit{学习一个简单的两层去噪自编码器}，\footnote{事实上，正如我们在上一章中提到的，PCA
        是神经网络最早被用来解决的问题之一
    \cite{Oja1982SimplifiedNM,Baldi89}。}这是非平凡神经网络最简单的例子之一。在这个框架中，\textit{学习到的表征}就是
    \((\vU^{\star})^{\top}\vx \approx \vz\)。通过这种方式，PCA
    充当了（深度）表征学习的模型问题，我们将在本专著中进一步借鉴。请注意，在这个类比中，表征反映了输入数据，或者是输入数据向学习到的低维结构的投影。这一属性在未来将特别相关。
\end{remark}

\subsection{通过幂迭代追求低秩结构}\label{subsec:power iterations}

有一种计算上高效的方法来估计 \(\vX\vX^{\top}/N\) 或任何对称半正定矩阵 \(\vM\)
的前几个特征向量，称为\textit{幂迭代 (power
iteration)}。该方法是我们将在本章后面讨论的几种高维数据分析算法方法的基础，因此我们在这里进行讨论。

设 \(\vM\) 为对称半正定矩阵。存在一个由 \(\vM\) 的特征向量 \((\vw_{i})_{i = 1}^{D}\) 组成的
\(\R^{D}\) 的正交基，其对应的特征值为 \(\lambda_{1} \geq \cdots \geq \lambda_{D}
\geq 0\)。根据定义，任何特征向量 $\vw_i$ 都满足 $\lambda_i \vw_i = \vM
\vw_i$。因此，对于任何 $\lambda_i > 0$，$\vw_i$ 都是以下方程的“不动点”：
\begin{equation}
    \vw = \frac{\vM\vw}{\norm{\vM\vw}_{2}}.
    \label{eqn:PCA-fixed-point}
\end{equation}

\begin{theorem}[幂迭代]
    假设对于所有 \(i > 1\)，都有 \(\lambda_{1} > \lambda_{i}\)。如果我们使用以下迭代计算
    \eqref{eqn:PCA-fixed-point} 的不动点：
    \begin{equation}\label{eq:power_iteration}
        \vv_{0} \sim \dNorm(\vzero, \vone), \qquad \vv_{t + 1} \gets
        \frac{\vM\vv_{t}}{\norm{\vM\vv_{t}}_{2}},
    \end{equation}
    那么，在极限情况下，\(\vv_{t}\) 将收敛到 \(\vM\) 的一个最大单位范数特征向量。
\end{theorem}

\begin{proof} 首先，请注意对于所有 \(t\)，我们有
    \begin{equation}
        \vv_{t} = \frac{\vM\vv_{t - 1}}{\norm{\vM\vv_{t - 1}}_{2}} =
        \frac{\vM^{2}\vv_{t - 2}}{\norm{\vM\vv_{t -
        1}}_{2}\norm{\vM\vv_{t - 2}}_{2}} = \cdots =
        \frac{\vM^{t}\vv_{0}}{\prod_{s = 1}^{t}\norm{\vM\vv_{s}}_{2}}.
    \end{equation}
    因此，\(\vv_{t}\) 与 \(\vM^{t}\vv_{0}\) 方向相同且为单位范数，所以我们可以写成
    \begin{equation}
        \vv_{t} = \frac{\vM^{t}\vv_{0}}{\norm{\vM^{t}\vv_{0}}_{2}}.
    \end{equation}
    因为 $\vM$ 的所有特征向量 \(\vw_{i}\) 构成了 \(\R^{D}\) 的正交基，我们可以写成
    \begin{equation}
        \vv_{0} = \sum_{i = 1}^{D}\alpha_{i}\vw_{i},
    \end{equation}
    其中，由于 \(\vv_{0}\) 是高斯分布的，\(\alpha_{i}\) 以概率 \(1\)
    全都不为零。因此，我们可以使用前面关于 \(\vv_{t}\) 的表达式来写成
    \begin{equation}
        \vv_{t} = \frac{\vM^{t}\vv_{0}}{\norm{\vM^{t}\vv_{0}}_{2}} =
        \frac{\sum_{i =
        1}^{D}\lambda_{i}^{t}\alpha_{i}\vw_{i}}{\norm{\sum_{i =
        1}^{D}\lambda_{i}^{t}\alpha_{i}\vw_{i}}_{2}} = \frac{\sum_{i
        = 1}^{D}\lambda_{i}^{t}\alpha_{i}\vw_{i}}{\sum_{i =
        1}^{D}\lambda_{i}^{t}\abs{\alpha_{i}}}.
    \end{equation}
    现在，让我们考虑 \(\lambda_{1} > \lambda_{2}
    \geq \cdots \geq \lambda_{D} \geq 0\) 的情况。（具有重复最大特征值的情况类似。）那么我们可以写成
    \begin{equation}
        \vv_{t} = \frac{\alpha_{1}\vw_{1} + \sum_{i =
        2}^{D}(\lambda_{i}/\lambda_{1})^{t}\alpha_{i}\vw_{i}}{\abs{\alpha_{1}}
        + \sum_{i = 2}^{D}(\lambda_{i}/\lambda_{1})^{t}\abs{\alpha_{i}}}.
    \end{equation}
    因为对于所有 \(i > 1\) 都有 \(\lambda_{1} > \lambda_{i}\)，求和号内的项以指数速度趋于 \(0\)，剩下的极限为
    \begin{equation}
        \lim_{t \to \infty}\vv_{t} =
        \frac{\alpha_{1}}{\abs{\alpha_{1}}}\vw_{1} = \sign(\alpha_{1})\vw_{1},
    \end{equation}
    这是 \(\vM\) 的一个最大单位特征向量。\(\vM\) 的最大特征值 \(\lambda_{1}\) 可以通过
    \(\vv_{t}^{\top}\vM\vv_{t}\) 来估计，它同样快速地收敛到 \(\lambda_{1}\)。
\end{proof}

\begin{remark}[作为去噪过程的幂迭代]
    请注意，通过幂迭代计算主成分可以被解释为一个“去噪”过程：从随机噪声开始，迭代收敛到一个一维子空间上。事实上，正如我们稍后将在
    \Cref{ch:denoising} 中展示的那样，如果我们假设{\em
    数据分布的支撑集是一个单一子空间}，相关的去噪算子（或得分函数）与上述幂迭代完全相同。在第 \ref{ch:denoising}
    章中，我们将展示如何将这种去噪过程推广到具有更一般低维结构的分布。这本质上是学习任何低维分布的一种通用且最有效的方法。
\end{remark}

为了找到第二个最大特征向量，我们将幂迭代算法应用于 \(\vM -
\lambda_{1}\vv_{1}\vv_{1}^{\top}\)，它具有特征向量 \((\vw_{i})_{i = 2}^{D}\)
和相应的特征值 \((\lambda_{i})_{i = 2}^{D}\)。通过依次重复此过程 \(d\)
次，我们可以非常高效地估计任何对称半正定矩阵 \(\vM\) 的前 \(d\) 个特征向量。因此，我们也可以将其应用于
\(\vX\vX^{\top}/N\) 以恢复前 \(d\)
个主成分，这正是我们最初的目标。请注意，这种方法一次恢复一个主成分；在未来的章节中，我们将把它与其他算法方法（例如对所有特征向量进行梯度下降）进行对比。

%\begin{remark}[Online PCA] For simplicity, all algorithms that we
% present in this chapter assume that a sufficient set of samples
% $\X$ have been given in advance. This is not always the case. In
% the case of continuous learning, we may need to learn the
% distribution (subspace) from samples that are given one at at time.
% This problem is often referred to as {\em online} PCA. We will
% leave the study to \Cref{ch:consistent-representations}.
%
%\end{remark}

\subsection{通过在线 PCA 进行局部与增量学习}
请注意，在上述构造中，用于编码和解码的线性变换
$\vU$（即主成分矩阵）是预先从所有输入数据中“离线”计算出来的。一个问题是，随着输入数据按顺序到来，这种变换是否可以“在线”学习。这个问题在
Oja 1982 年的工作 \cite{Oja1982SimplifiedNM} 中得到了解答。
\begin{example}[PCA 的归一化赫布学习方案] 考虑一个独立同分布的随机向量序列 $\x_1, \ldots, \x_i, \ldots \in
    \mathbb{R}^n$，其协方差为 $\boldsymbol{\Sigma} \in
    \mathbb{R}^{n\times n}$。设 $\vu_0 \in \mathbb{R}^n$，并将输入向量 $\x_i$
    对权重向量 $\vu_i$ 的响应定义为它们的内积：
    \begin{equation}
        \eta_i = \vu_i^T \x_i
    \end{equation}
    我们根据以下方案更新权重向量：
    \begin{equation}
        \vu_{i+1} = \frac{\vu_i + \gamma \eta_i \x_i}{\|\vu_i + \gamma
        \eta_i \x_i\|_2}
        \label{eqn:Hebbian}
    \end{equation}
    其中 $\gamma >0$ 为某个较小的增益。这种更新方案可以被视为一种归一化的赫布 (Hebbian) 方案，其中如果输入
    $\x$ 和输出 $\eta$（的乘积）都很强，神经元之间连接的权重就会变得更强。人们可以将权重向量 $\vu$ 视为基于来自输出
    $\eta$ 的某种形式的反馈而“学习”到的。然后，在合理的假设下，Oja \cite{Oja1982SimplifiedNM}
    证明了 $\vu_i$ 收敛到与 $\boldsymbol{\Sigma}$ 的最大特征值相关的特征向量。
\end{example}

为了解释归一化赫布方案 \eqref{eqn:Hebbian}，让我们将数据样本 PCA 目标
\eqref{eq:orthogonal_denoising} 视为目标\textit{总体}版本的一个近似，该总体版本通过对样本
$\vx$ 分布的期望来定义：
\begin{equation}\label{eq:pca-reconstruction-expectation}
    \min_{\vU^\top \vU = \vI}\, \bE_{\vx}\left[\norm*{\vx
    - \vU\vU^\top \vx}_2^2\right]
\end{equation}
我们将在下一节进一步深入探讨这个目标。那么，归一化赫布方案可以被解释为对
\Cref{eq:pca-reconstruction-expectation}
的\textit{随机}投影梯度下降方案的一阶近似（批量大小为 $1$，且 $\vU$ 的列数为 $1$），只要
$\norm{\vu}_2 = 1$，这由 \eqref{eqn:Hebbian}
中的投影操作来维持。这种\textit{在线算法}是处理具有因果结构的顺序结构数据的更复杂算法的有用原语，正如我们将在
\Cref{sec:causal-crate} 和 \Cref{sec:clm_text}
中进一步研究的那样。有趣的是，更一般地说，简单的在线 PCA
算法暗示了这样一个事实：\textit{比（端到端）反向传播更简单的算法也能成功学习编码器-解码器对}，我们将在
\Cref{ch:future} 中简要回顾这一点，作为未来智能研究的一个重要目标。

\subsection{统计视角：概率 PCA}\label{subsec:probabilistic PCA}

请注意，上述公式没有对数据生成过程做出任何统计假设。然而，在给定的数据模型中包含统计元素是很常见的，因为它可能会对分析结果增加进一步的启发性解释。因此，我们提出了一个自然的问题：\textit{低维结构的统计类比是什么？}我们的回答是，低维\textit{分布}是指其支撑集集中在低维几何结构周围的分布。

为了说明这一点，我们讨论\textit{概率主成分分析 (PPCA)} \citep{TippingM1999}。这种表述可以被视为常规
PCA 的统计变体。在数学上，我们现在将数据视为取值于 \(\R^{D}\) 的随机变量
\(\vx\)（有时也称为\textit{随机向量}）的样本。当且仅当存在一个正交矩阵 \(\vU \in \O(D,
d)\)、一个取值于 \(\R^{d}\) 的随机变量 \(\vz\) 以及一个取值于 \(\R^{D}\)
的\textit{较小}随机变量 \(\veps\)，使得 \(\vz\) 和 \(\veps\) 独立，并且满足以下条件时，我们称
\(\vx\) 具有（近似）低秩统计结构：
\begin{equation}
    \vx = \vU\vz + \veps.
\end{equation}
我们的目标仍然是恢复 \(\vU\)。为此，我们建立了与 \Cref{sub:pca} 中类似的问题，即对子空间支撑
\(\tilde{\vU}\) 和随机变量 \(\vz\) 进行优化，以解决以下问题：
\begin{equation}
    \min_{\tilde{\vU}, \tilde{\vz}}\Ex\norm{\vx -
    \tilde{\vU}\tilde{\vz}}_{2}^{2}.
\end{equation}
由于我们正在寻找最佳的此类随机变量 \(\vz\)，我们可以针对 \(\vx\) 的每个值分别找到其实现 \(\vz(\vx)\)。执行与
\Cref{sub:pca} 中相同的计算，我们得到 %
\begin{equation}\label{eq:ppca_denoising}
    \min_{\tilde{\vU}, \tilde{\vz}}\Ex\norm{\vx -
    \tilde{\vU}\tilde{\vz}}_{2}^{2} =
    \min_{\tilde{\vU}}\Ex\min_{\tilde{\vz}(\vx)}\norm{\vx -
    \tilde{\vU}\tilde{\vz}(\vx)}_{2}^{2} =
    \min_{\tilde{\vU}}\Ex\norm{\vx - \tilde{\vU}\tilde{\vU}^{\top}\vx}_{2}^{2},
\end{equation}
再次强调了这样一个事实：具有主成分 \(\vU^{\star}\) 的估计子空间对应于一个投影到该子空间上的去噪器
\(\vU^{\star}(\vU^{\star})^{\top}\)。和以前一样，我们得到
\begin{align}\label{eq:PPCA}
    \argmin_{\tilde{\vU}}\Ex\norm{\vx -
    \tilde{\vU}\tilde{\vU}^{\top}\vx}_{2}^{2}
    &= \argmax_{\tilde{\vU}}\Ex\norm{\tilde{\vU}^{\top}\vx}_{2}^{2} \\
    &=
    \argmax_{\tilde{\vU}}\tr(\tilde{\vU}^{\top}\Ex[\vx\vx^{\top}]\tilde{\vU}),
\end{align}
而后一个问题的解正是二阶矩矩阵 \(\Ex[\vx\vx^{\top}]\) 的前 \(d\)
个主成分。实际上，上述问题在视觉上与上一小节中计算主成分的方程非常相似，只是用 \(\Ex[\vx\vx^{\top}]\) 替换了
\(\vX\vX^{\top}/N\)。事实上，后者是前者的一个估计。这两种表述实际上做的是同一件事，并且具有相同的实际解决方案——计算数据矩阵
\(\vX\) 的左奇异向量，或者等价地，计算估计协方差矩阵 \(\vX\vX^{\top}/N\)
的前几个特征向量。然而，统计表述有一个额外的解释。假设 \(\Ex[\vz] = \vzero\) 且 \(\Ex[\veps] =
\vzero\)。我们有
\begin{equation}
    \Ex[\vx] = \vU\Ex[\vz] + \Ex[\veps] = \vzero,
\end{equation}
因此 \(\Cov[\vx] = \Ex[\vx\vx^{\top}]\)。现在计算 \(\Cov[\vx]\)，我们有
\begin{equation}
    \Cov[\vx] = \vU\Cov[\vz]\vU^{\top} + \Cov[\veps] =
    \vU\Ex[\vz\vz^{\top}]\vU^{\top} + \Cov[\veps].
\end{equation}
特别是，如果 \(\Cov[\veps]\) 很小，那么 \(\Cov[\vx]
= \Ex[\vx\vx^{\top}]\) 近似为一个低秩矩阵，特别是秩为
\(d\)。因此，\(\Ex[\vx\vx^{\top}]\) 的前 \(d\)
个特征向量本质上总结了整个协方差矩阵。但它们也是主成分，因此我们可以将主成分分析解释为对 \(\Cov[\vx]\) 执行低秩分解。

\begin{remark}
    通过使用 PCA 的概率视角，我们对其与去噪的关系有了更清晰、更定量的理解。首先，考虑
    \eqref{eq:ppca_denoising} 中的去噪问题，即
    \begin{equation}
        \min_{\tilde{\vU}}\Ex\norm{\vx -
        \tilde{\vU}\tilde{\vU}^{\top}\vx}_{2}^{2}.
    \end{equation}
    不难证明，如果 $\tilde{\vU}$ 有 $d$ 列，并且如果 \(\veps\) 是各向同性的高斯随机变量，即分布为
    \(\veps \sim \dNorm(\vzero,
    \sigma^{2}\vI)\)，\footnote{只要支撑整个
    \(\R^{D}\)，其他分布也可以，但高斯分布在这里最容易处理。}那么对于该问题的\textit{任何}最优解 \(\vU^{\star}\)，我们有
    \begin{equation}
        \vU^{\star}(\vU^{\star})^{\top} = \vU\vU^{\top}
    \end{equation}
    因此，真实的支撑子空间（例如 \(\cS \doteq
    \mathop{\mathrm{col}}(\vU)\)）被恢复为 \(\vU^{\star}\) 的列的张成空间，因为
    \begin{equation}
        \cS = \mathop{\mathrm{col}}(\vU) = \mathop{\mathrm{col}}(\vU\vU^{\top})
        = \mathop{\mathrm{col}}(\vU^{\star}(\vU^{\star})^{\top})
        = \mathop{\mathrm{col}}(\vU^{\star}).
    \end{equation}
    特别是，学习到的\textit{去噪映射} \(\vU^{\star}(\vU^{\star})^{\top}\) 是到
    \(\cS\) 上的正交投影，将嘈杂点推向真实的支撑子空间。在只有有限样本的情况下（如 \Cref{thm:pca}
    中），我们可以建立类似的技术结果，但这需要更多的努力和技术性。总结这一讨论，我们得到以下非正式定理。
\end{remark}

\begin{theorem}\label{thm:ppca}
    假设随机变量 \(\vx\) 可以写成
    \begin{equation}
        \vx = \vU\vz + \veps
    \end{equation}
    其中 \(\vU \in \O(D, d)\) 捕获了\textit{低秩结构}，\(\vz\) 是取值于 \(\R^{d}\)
    的随机向量，\(\veps\) 是取值于 \(\R^{D}\) 的随机向量，使得 \(\vz\) 和 \(\veps\) 独立，且
    \(\veps\) 较小。那么我们数据集的\textit{主成分} \(\vU^{\star} \in \O(D, d)\) 由
    \(\Ex[\vx\vx^{\top}]\) 的前 \(d\)
    个特征向量给出，并且近似对应于最优线性去噪器：\(\vU^{\star}(\vU^{\star})^{\top}
    \approx \vU\vU^{\top}\)。
\end{theorem}

\subsection{掩码低秩矩阵补全}

在前面的小节中，我们讨论了\textit{学习低秩几何或统计分布}的问题，其中数据是从带有加性噪声的子空间中采样的。但这并不是唯一值得研究的来自低维分布的扰动类型。在本小节中，我们介绍另一类在深度学习中变得越来越重要的非加性误差。让我们考虑这样一种情况：我们有一些根据
\eqref{eq:pca_dgp} 生成的数据 \(\{\vx_{i}\}_{i = 1}^{N}\)。现在我们将它们排列成一个矩阵
\(\vX = \mat{\vx_{1}, \dots, \vx_{N}} \in \R^{D \times
N}\)。与以前不同，我们不直接观察 \(\vX\)；相反，我们假设我们的观测在途中被破坏，我们得到的是
\begin{equation}
    \vY = \vM \hada \vX,
\end{equation}
其中 \(\vM \in \{0, 1\}^{D \times N}\) 是我们已知的\textit{掩码
(mask)}，\(\hada\) 是逐元素乘法。在这种情况下，我们的目标是仅给定损坏的观测值 \(\vY\)、掩码 \(\vM\) 以及
\(\vX\) 是（近似）低秩的知识，来恢复 \(\vX\)（从这一点开始，我们可以使用 PCA 来恢复 \(\vU\)
等）。这被称为\textit{低秩矩阵补全}。

这个以及这种低秩结构恢复问题的类似推广被称为“鲁棒 PCA (Robust
PCA)”问题。已经有许多优秀的资源讨论了解决这些问题的算法和方法
\cite{Wright-Ma-2022}。因此，我们在这里不深入探讨解决方法。相反，我们将讨论在什么条件下解决这个问题是\textit{合理的}。一方面，在最荒谬的情况下，假设矩阵
\(\vX\) 的每个元素都是独立于所有其他元素选择的。那么，即使只缺失一个元素而我们有 \(DN - 1\) 个元素，也没有希望完全恢复
\(\vX\)。另一方面，假设我们知道 \(\vX\) 的秩恰好为 1，这是对数据低维结构的一个极其强烈的条件，并且我们正在处理掩码
\begin{equation}
    \vM = \mat{\vone_{(D - 1) \times 1} & \vzero_{(D - 1) \times (N -
    1)} \\ 1 & \vone_{1 \times (N - 1)} }.
\end{equation}
那么我们知道数据分布在一条直线上，并且我们知道该直线上的一条向量——它就是矩阵 \(\vY = \vM
\hada \vX\)
的第一列。从每一列的最后一个坐标（也由掩码向我们揭示），我们可以求解出每一列，因为对于每个最终坐标，直线上只有一个向量具有该坐标。因此，我们可以以完美的精度重建
\(\vX\)，并且我们只需要线性数量的观测值 \(D + N - 1\)。

在现实世界中，实际问题介于上述两种极端情况之间。然而，这两个极端之间的差异，以及前面关于 PCA 的讨论，揭示了一个普遍的真理核心：
\begin{quote}
    \centering
    \textit{数据分布的维度越低、结构越强，就越容易处理，需要的观测值也就越少——前提是算法有效地利用了这种低维结构。}
\end{quote}
也许可以预见，在手稿的其余部分，从下一节开始，我们将反复遇到这个主题。\footnote{在后面的章节中，我们将看到，人们可以仅用一小部分未掩码的观测值来补全现实世界的数据（如图像），例如
\Cref{sec:image_completion} 中图像的掩码自编码。这也是因为自然图像的分布在所有图像的空间中具有相当低维（尽管是非线性）的结构。}

\section{完备低维子空间的混合}%
\label{sec:ica}
正如我们所见，低秩信号模型足够丰富，能够全面展示数据低维性与用于表征和在误差下恢复的高效、可扩展计算算法之间的相互作用。这些模型暗示了一个\textit{线性}且对称的表征学习流水线
\eqref{eqn:autoencoding-PCA}：
\begin{equation*}
    \vz = \cE(\vx) = \vU^\top \vx, \quad \hat{\vx} = \cD(\vz) = \vU \vz,
\end{equation*}
只要 $\vx$ 的分布确实是线性的，就可以通过主成分分析（例如，使用幂方法高效求解）从 $\vx$
的有限样本中可证明地学习到该流水线。这是一个限制性的假设——正如 20 世纪杰出的统计学家哈罗德·霍特林 (Harold
Hotelling)\footnote{巧合的是，他也因对主成分分析的发展和命名做出的贡献而闻名
\cite{Hotelling1933}。}在乔治·丹齐格 (George Dantzig) 首次提出其线性规划理论后所反驳的那样
\cite{Dantzig2002-eh}，
\begin{quote}
    \centering
    \textit{“……我们都知道世界是非线性的。”}
\end{quote}

\begin{figure}
    \centering
    \includegraphics[height=0.35\linewidth]{\toplevelprefix/chapters/chapter2/figs/motion.png}
    \hspace{5mm}
    \includegraphics[height=0.349\linewidth]{\toplevelprefix/chapters/chapter2/figs/segment.png}
    \caption{\textbf{左：}在场景中独立移动的物体上跟踪的特征。\textbf{右：}与图像不同区域相关的图像块。}
    \label{fig:multiple-subspaces}
\end{figure}
即使考虑到其优雅和简单，低秩假设也过于严格，无法广泛应用于现实世界数据的建模。一个关键的局限性是假设由\textit{单一}线性子空间负责生成结构化观测值。在许多实际应用中，由不同低维子空间的\textit{混合}生成的结构提供了一个更现实的模型。例如，考虑一个捕获几个不同物体运动的视频序列，每个物体都有其独立的位移（\Cref{fig:multiple-subspaces}
左）。在对个体运动的适当假设下，每个物体在拼接的视频帧序列中负责一个独立的低维子空间
\cite{VidalR2004-ECCV}。再举一个例子，考虑通过学习图像内\textit{图像块
(patches)}（空间上连续的像素集合）分布的模型来对自然图像进行建模（\Cref{fig:multiple-subspaces}
右）。与 \Cref{fig:faces-digits}
中简单的设置（其中姿态匹配的人脸图像可以被单个低维子空间很好地近似）不同，自然图像中特定位置的图像块可以对应于具有非常不同属性的物体——例如，由于遮挡边界而导致的明显颜色或形状。因此，用单个子空间对图像块的分布进行建模是徒劳的，但是子空间的\textit{混合}（每个区域一个）在实践中表现得非常好，例如用于分割或压缩目的
\cite{Mobahi-IJCV2011}。\footnote{我们将在 \Cref{ch:deep-networks}
中回到这一观察结果，在那里我们将展示它可以被显著推广以学习大规模现代数据集的表征。}我们将在后面的章节（\Cref{eg:image-segmentation}）中看到一个具体的例子。

% \begin{figure}
%     \centering
%     \includegraphics[height=0.35\linewidth]{\toplevelprefix/chapters/chapter2/figs/subspaces.png}
%     \caption{Data on a mixture of low-dimensional subspaces, say
% $\mathcal{S}_j
%     = \mathop{\mathrm{col}}(\vU_j)$.}
%     \label{fig:subspaces}
% \end{figure}

在本节中，我们将首先讨论当数据分布可以由{\em 低维子空间的混合}建模时，结构化表征学习的概念和算法基础，如
\Cref{fig:subspaces} 所示。在这种设置下，解码器映射将几乎与单个子空间的情况一样简单：我们只需通过以下方式表示
\begin{equation}\label{eq:mixture-subspaces-decoder-first-cut}
    \hat{\x} = \cD(\vz) = \left( \sum_{k=1}^K \pi_k(\vz)\vU_k \right) \vz,
\end{equation}
其中 $\pi_k : \R^d \to \set{0,1}$
是一组\textit{稀疏}加权随机变量，使得在求和中只选择一个单一的子空间 $\mathcal{S}_k
= \mathop{\mathrm{col}}(\vU_k)$。然而，将此类数据 $\vx$ 编码为合适的表征
$\vz$，并从数据中学习这样的编码器-解码器对的任务，将被证明更加复杂。

我们将看到，来自\textit{稀疏表征}和\textit{独立成分分析}丰富文献中的思想，如何通过稀疏性的视角自然地重新表述上述解码器架构，以及相应的编码器架构（通过类似于主成分分析的幂方法算法获得），并为从数据中学习此类编码器-解码器对提供正确性和效率的强有力保证。从这个意义上说，混合线性低维结构的情况已经引出了结构化表征学习的许多关键概念组件，我们将在本书中以更广泛的一般性来发展这些组件。

\subsection{子空间混合与稀疏使用字典}\label{sec:mixture-and-dict}

设 $\vU_1, \dots, \vU_K$（每个大小为 $D \times d$）表示 $\R^D$ 中 $K$ 个维度为 $d$
的子空间的正交基集合。从几何上讲，说 $\vx$ 服从由 $\vU_1, \dots, \vU_K$ 参数化的子空间混合分布，意味着
\begin{equation}\label{eq:mixture-of-subspaces-geometric}
    \vx = \vU_k \vz  \quad \text{for some} \enspace k \in
    [K],\enspace \vz \in \R^d.
\end{equation}
正如我们在 PCA 和线性结构的情况下所看到的，这种几何模型的统计类比是 $\vx$ 服从\textit{高斯混合}分布：即，
\begin{equation}\label{eq:mixture-of-subspaces-statistical}
    \vx \sim \sum_{k=1}^K \pi_k \cN(\mathbf{0}, \vU_k\vU_k^\top),
    \quad \text{for some} \enspace \pi_k \geq 0,\enspace \sum_{k=1}^K \pi_k = 1.
\end{equation}
换句话说，对于每个 $k \in [K]$，$\vx$ 以概率 $\pi_k$ 在低维子空间
$\mathop{\mathrm{col}}(\vU_k)$ 上呈高斯分布。

\begin{remark}[高斯混合与高斯叠加]
    应该注意的是，上述模型 \eqref{eq:mixture-of-subspaces-statistical}
    是高斯分布的\textit{混合}，不要与通过叠加形成的高斯变量混合相混淆，例如
    \begin{equation}
        \vx = \sum_{i=1}^n w_i \vx_i, \quad \vx_i \sim
        \cN(\mathbf{0}, \vU_i\vU_i^\top),
    \end{equation}
    其中 $\vx_i$ 是独立的随机高斯向量，$w_i$ 是一组固定权重。正如我们从高斯向量的性质中所知，这种叠加 $\vx$ 仍将是高斯分布。
\end{remark}

现在，我们关注 \eqref{eq:mixture-of-subspaces-geometric}
提供的几何视角。对于这种条件表示，有一种在代数上很方便的替代方法。考虑一个\textit{提升的 (lifted)} 表征向量 $\vz
= [\vz_1^\top, \dots, \vz_K^\top]^\top \in \R^{dK}$，使得 $\vz$
是\textit{$d$-稀疏的}，其支撑集位于 $dK$ 个坐标中 $K$ 个连续且不重叠的 $d$ 坐标块之一上。那么
\eqref{eq:mixture-of-subspaces-geometric} 可以等价地写为
\begin{equation}\label{eq:mixture-of-subspaces-dictionary-pre}
    \vx =
    \underbrace{
        \begin{bmatrix}
            | & \dots & |  \\
            \vU_1 & \dots & \vU_K  \\
            | & \dots & |
        \end{bmatrix}
    }_{\vU}
    \underbrace{
        \begin{bmatrix} \vz_1 \\ \vdots \\ \vz_K
        \end{bmatrix}
    }_{\vz},
    \quad
    \norm*{
        \begin{bmatrix} \norm*{\vz_1}_2 \\ \vdots \\ \norm*{\vz_K}_2
        \end{bmatrix}
    }_0 = 1.
\end{equation}
这里，$\ell^0$ “范数” $\norm{\,\cdot\,}_0$ 通过计算非零元素的数量来衡量稀疏性：
\begin{equation}\label{eq:ell-zero-norm}
    \norm{\vz}_0 = \abs*{\set{i \given z_i \neq 0}},
\end{equation}
矩阵 $\vU \in \R^{D \times Kd}$ 被称为\textit{字典}，所有 $\{\vU_i\}_{i=1}^K$
作为码字。一般来说，如果混合中的子空间数量 $K$ 足够大，那么字典 $\vU$ 中包含的列数就没有上限。在 $Kd < D$
的情况下，$\vU$ 被称为\textit{欠完备的 (undercomplete)}；当 $Kd = D$
时，它被称为\textit{完备的 (complete)}；当 $Kd > D$ 时，它被称为\textit{过完备的 (overcomplete)}。

现在，\eqref{eq:mixture-of-subspaces-dictionary-pre}
提出了一种为了分析的易处理性而进行的便捷松弛：与其将 $\vx$ 建模为来自 $K$ 个\textit{特定}子空间 $\vU_1,
\dots, \vU_K$ 的混合，我们不如从一个字典 $\vU \in \R^{D \times m}$ 开始（其中 $m$
可以小于或大于 $D$），并简单地寻求用足够小的稀疏度 $\norm{\vz}_0$ 来表示 $\vx = \vU \vz$。这引出了
$\vx$ 的\textit{稀疏字典模型}：
\begin{equation}\label{eq:mixture-of-subspaces-dictionary}
    \vx =  \vU \vz + \veps,
    \quad
    \norm{\vz}_0 \ll d,
\end{equation}
其中 $\veps$ 表示未知的噪声向量。从几何上讲，这意味着 $\vx$ 靠近 $\vU$ 的 $\norm{\vz}_0$
列子集的张成空间，这使其成为混合子空间模型 \eqref{eq:mixture-of-subspaces-geometric}
的一个实例，具有非常大的 $K$ 值，以及子空间 $\vU_k$ 之间的特定相关性。

\paragraph{用于稀疏编码的正交字典。}
现在我们可以表述本节将要研究的低维子空间混合的结构化表征学习问题。我们假设我们有来自未知稀疏字典模型
\eqref{eq:mixture-of-subspaces-dictionary} 的样本 $\vX = [\vx_1, \dots,
\vx_N]$，可能带有加性噪声 $\veps_i$。让我们从这样一个假设开始：稀疏字典模型
\eqref{eq:mixture-of-subspaces-dictionary} 中的字典 $\vU$
是完备且正交的，\footnote{可以证明，对于完备的情况，做出正交假设并不会失去任何一般性（\Cref{exercise:whitening}）。}并且每个系数向量
$\vz$ 都是 $d$-稀疏的，其中 $d \ll D$。
```latex
In this setting, representation learning amounts to correctly
learning the orthogonal dictionary $\vU$ via optimization: we
can then take $\cE(\vx) = \vU^\top \vx$ as the encoder and $\cD(\vz) = \vU \vz$
as the decoder, and $\cD = \cE^{-1}$. In diagram form:
\begin{equation}
    \vx \xrightarrow{\hspace{2mm} \mathcal{E} = \vU^\top
    \hspace{2mm}}  \vz \xrightarrow{\hspace{2mm} \mathcal{D} = \vU
    \hspace{2mm}}   \hat{\vx}.
    \label{eqn:autoencoding-DL}
\end{equation}
我们看到，完备字典学习的自编码对 $(\cE, \cD)$
是对称的，就像单一线性子空间的情况一样，这使得编码和解码的计算任务并不比线性情况更困难。另一方面，学习字典 $\vU$
的任务严格来说比通过 PCA 学习单一线性子空间更困难。
要理解为什么我们不能简单地使用 PCA 来正确学习正交字典 $\vU$，请注意产生 PCA 的损失函数，即
\eqref{eq:pca_equals_denoising}，对矩阵 $\vU$ 的行的旋转是完全不变的：也就是说，如果 $\vQ$
是任意 $d \times d$ 正交矩阵，那么 $\vU$ 和 $\vU \vQ$ 都是可行的，并且对于
\eqref{eq:pca_equals_denoising} 具有相同的损失。稀疏字典模型显然对这种变换不是不变的：如果我们用 $\vU
\vQ$ 替换 $\vU$，并对表征系数 $\vz$ 进行相应的旋转 $\vQ^\top \vz$，我们通常会破坏 $\vz$
的稀疏结构，从而违反建模假设。因此，我们需要用于学习正交字典的新算法。

\subsection{完备字典学习}
\label{sec:complete-dictionary}

在本节中，我们将推导解决正交字典学习问题的算法。更准确地说，我们假设观测向量 $\vx \in \R^D$ 遵循统计模型
\begin{equation}
    \vx = \vU \vz + \veps,
    \label{eq:ica-model-ch2}
\end{equation}
其中 $\vU \in \R^{D \times D}$ 是一个未知的正交字典，$\vz$ 是一个具有统计独立分量 $z_i$
的随机向量，每个分量均值为零，而 $\veps \in \R^D$ 是一个由小（高斯）噪声组成的独立随机向量。目标是从 $\vx$
的样本中恢复 $\vU$（从而恢复 $\vz$）。

这里我们假设每个独立分量 $z_i$ 的分布为 $$z_i \sim \dBern(\theta) \cdot \cN(0,
1/\theta).$$ 也就是说，它是一个取值为 $1$ 的概率为 $\theta$、取值为 $0$ 的概率为 $1-\theta$
的伯努利随机变量，与一个方差为 $1/\theta$ 的独立高斯随机变量的乘积。这种分布正式称为 {\em
伯努利-高斯}（Bernoulli-Gaussian）分布。选择这种归一化是为了使 $\Var(z_i) = 1$，从而
$\bE[\norm{\vz}_2^2]=d$。这种建模假设意味着独立分量向量 $\vz$ 通常非常稀疏：我们计算出
$\bE\left[\norm{\vz}_0\right] = d\theta$，当 $\theta$ 与 $d$ 成反比时，这个值很小。

\begin{remark}[正交假设]
    乍一看，字典 $\vU$ 是正交的这一假设似乎有些局限。但实际上这并不失一般性。我们可以将完备字典视为任何方形可逆矩阵
    $\vU$。对于由该字典生成的样本：$\vX = \vU \vZ \in
    \mathbb{R}^{D\times N}$，很容易证明，通过对数据矩阵 $\vX$ 进行一些预处理：
    \begin{equation}
        \bar{\vX} = \Big(\frac{1}{N\theta} \vX\vX^\top\Big)^{-\frac{1}{2}}\vX,
    \end{equation}
    那么存在一个正交矩阵 $\vU_{o} \in \O(D)$ 使得
    \begin{equation}
        \bar{\vX} = \vU_{o}\vZ.
    \end{equation}
    更多细节请参见 \Cref{exercise:whitening} 或 \cite{sun2017completeI}。
\end{remark}

\begin{figure}
    \centering
    \includegraphics[width=0.6\linewidth]{\toplevelprefix/chapters/chapter2/figs/2DL4Sphere.png}\vspace{-0.1in}
    \caption{最大化 $\ell^4$ 范数或最小化 $\ell^1$ 范数可以促进稀疏性（对于球面上的向量）。}
    \label{fig:L4-sphere}
\end{figure}

\paragraph{通过 MSP 算法进行字典学习。}

现在假设我们给定一组观测值：
\begin{equation}
    \vx_i = \vU \vz_i + \veps_i,\ \forall i \in [N].
\end{equation}
令 $\vX = [\vx_1, \dots, \vx_N]$ 且 $\vZ = [\vz_1, \dots,
\vz_N]$。目标是从数据 $\vX$ 中恢复 $\vU$。因此，给定任意正交矩阵 $\vA \in \O(D)$，
\begin{equation}
    \vA\vx_i = \vA\vU \vz_i + \vA\veps_i
\end{equation}
如果 $\vA = \vU^T$，它将几乎是稀疏的（因为根据假设，噪声 $\veps_i$ 的幅度很小）。

此外，鉴于 $\vU$ 是正交的且 $\veps$ 很小，向量 $\vx$ 具有可预测的期望范数，即
$\bE[\norm{\vx}_2^2] \approx
\bE[\norm{\vz}_2^2]=d$。一个已知的事实是，对于球面上的向量，最大化 $\ell^4$ 范数等价于最小化
$\ell^0$ 范数（以促进稀疏性），
\begin{equation}
    \argmax_{\vz \in
    \mathbb{S}^n}\|\vz\|_{4}\quad\Leftrightarrow\quad \argmin_{\vz
    \in \mathbb{S}^n}\|\vz\|_{0}.
\end{equation}
这在 \Cref{fig:L4-sphere} 中进行了说明。

正交矩阵 $\vA$ 保持欧几里得（\(\ell^{2}\)）范数不变：$\|\vA \vx\|_2^2 =
\|\vx\|_2^2$。因此，为了从 $\vX$ 中找到正确的正交字典 $\vU$，我们可以尝试求解以下优化程序：
\begin{equation}\label{eq0:orthogonal-dictionary-learning-l4}
    \max_{\tilde{\vA} \in \O(D)}\,
    \frac{1}{4} \norm*{
        \tilde{\vA} \vX
    }_4^4 =  \frac{1}{4} \sum_{i=1}^N \norm*{
        \tilde{\vA} \vx_i
    }_4^4
\end{equation}
这被称为 $\ell^4$ 最大化问题 \cite{Zhai-2020}。在找到解 \(\vA^{\star}\) 之后，我们可以取其转置
\(\vU^{\star}
= (\vA^{\star})^{\top}\)。
\begin{remark}
    众所周知，对于球面上的向量，最小化 $\ell^1$ 范数等价于最小化 $\ell^0$ 范数（以促进稀疏性），
    \begin{equation*}
        \argmin_{\vz \in
        \mathbb{S}^n}\|\vz\|_{1}\quad\Leftrightarrow\quad
        \argmin_{\vz \in \mathbb{S}^n}\|\vz\|_{0},
    \end{equation*}
    这同样在 \Cref{fig:L4-sphere} 中进行了说明。这一事实也可以被用来有效且高效地学习字典
    $\vA$。实际上，这比此处使用的 $\ell^4$ 范数探索得更早。感兴趣的读者可以参考
    \cite{qu2020findingsparsestvectorssubspace} 的工作。
\end{remark}

请注意，上述问题等价于以下约束优化问题：
\begin{equation}\label{eq:orthogonal-dictionary-learning-l4}
    \min\,
    -   \frac{1}{4} \norm*{
        \tilde{\vA} \vX
    }_4^4 \quad \mbox{受限于} \quad  \tilde{\vA}^\top\tilde{\vA} = \vI.
\end{equation}
如 \cite{Wright-Ma-2022} 所示，使用拉格朗日乘子法，可以推导出该问题的最优解应满足以下“不动点”条件：
\begin{equation}
    \vA^{\star} = \mathcal{P}_{\mathrm O(D)}[( {\vA^{\star}
    \vX})^{\hada 3}\vX^\top],
\end{equation}
其中 $\mathcal{P}_{\mathrm O(D)}[\,\cdot\,]$ 是到正交矩阵空间 $\mathrm O(D)$
的投影。\footnote{对于任意具有 SVD $\vM = \vU\bm \Sigma \vV^\top$ 的矩阵
    $\vM$，$\mathcal{P}_{\mathrm
O(D)}[\vM] = \vU \vV^\top$。我们将此作为练习留给读者。}

为了计算上述方程的不动点，类似于我们计算 PCA 特征向量 \eqref{eqn:PCA-fixed-point} 的方式，我们可以采用以下幂迭代：
\begin{equation}\label{eq:msp_iteration}
    \vA_{t+1} = \mathcal{P}_{\mathrm O(D)}[( {\vA_t \vX})^{\hada 3}\vX^\top].
\end{equation}
这被称为由 \cite{Zhai-2020} 提出的 {\em 匹配、拉伸和投影}（matching, stretching, and
projection，MSP）算法。研究表明，在宽泛的条件下，这种贪心算法确实能以超线性收敛率收敛到正确的解。

\begin{remark}[$\ell^4$ 最大化的全局最优性]\label{rem:L4-global}
    带约束的 $\ell^4$
    最大化问题是一个非凸规划。通常情况下，人们\textit{不应}期望任何贪心算法（例如梯度下降类型的算法）会收敛到全局最优解。令人惊讶的是，可以证明，与一般的非凸规划不同，球面上的
    $\ell^4$ 最大化
    \begin{equation}\label{eq:l4-maximization-sphere}
        \min\,
        -   \frac{1}{4} \norm*{
            \vq^\top \vX
        }_4^4 \quad \mbox{受限于} \quad  \vq^\top\vq = 1.
    \end{equation}
    的优化景观非常良好：所有局部极小值都接近全局最优值，并且所有临界点都是具有负曲率方向的鞍点。因此，任何具有逃离严格鞍点能力的下降方法都可以被证明能找到全局最优解。有关更精确的陈述，感兴趣的读者可以参考
    \cite{Qu2020Geometric}。
\end{remark}

\begin{remark}[稳定的深度线性网络]
    上述计算字典的迭代过程具有一种自然的增量式“深度学习”解释。让我们定义 $\delta \vA_{t+1} =
    \vA_{t+1}\vA_{t}^\top$ 和 $\vZ_t = \vA_t
    \vX$，那么很容易证明
    $$\delta \vA_{t+1} = \mathcal{P}_{\mathrm O(D)}[(\vZ_t)^{\hada 3}
    \vZ_t^\top].$$
    如果 $\vA_t$ 收敛到正确的字典 $\vD_o$，那么上述迭代编码过程本质上等价于一个“深度线性网络”：
    $$\vZ \; \longleftarrow \; \vZ_{t+1} =  \underbrace{\delta
    \vA_{t+1} \delta \vA_{t} \ldots \delta \vA_{1}}_{\color{red}
    \text{前向构建的层}} \vX.$$
    请注意，每一层增量变换 $\delta \vA_{t+1}$ 的计算仅依赖于前一层 $\vZ_t$
    的特征输出。该网络天然是稳定的，因为每一层都是保范数的正交变换。尽管它类似于线性深度网络，但学习每一层并不需要反向传播。所有层都在一次前向传递中学习完毕！
\end{remark}

\subsection{与 ICA 和峰度的联系}
在伯努利-高斯模型下，变量 $z_i$
是独立且非高斯的。那么，字典学习与经典的独立成分分析（ICA）之间存在着明确的对应关系，以至于解决其中一个问题的算法可以用来解决另一个问题。\footnote{我们在
    \Cref{exercise:symmetry-identifiability}
    中更深入地探讨了这个问题，其中建立了独立分量的非高斯性与纯几何对称性概念之间的联系。这个问题与我们上面的观察有关，即 PCA
无法用于恢复稀疏使用的正交字典：在统计设定中，它可以与高斯分布的旋转不变性联系起来（\Cref{exercise:gaussian-rot-invar}）。}

为了推导基于 ICA 的算法，我们关注一个称为\textit{峰度}（kurtosis）的目标函数，它在 ICA
中的使用是分量非高斯性的直接结果。零均值随机变量 $X$ 的\textit{峰度}（或四阶累积量）定义为
\begin{equation}\label{eq:kurtosis}
    \kurt(X) = \Ex{X^4} - 3 (\Ex{X^2})^2.
\end{equation}
如果我们只有来自随机变量 $X$ 的有限样本，并将其排列成向量 $\vx = [x_1, \dots, x_N]$，我们通过它们的经验平均值来定义峰度，这得出
\begin{equation}\label{eq:kurtosis-vector}
    \kurt(\vx) = \frac{1}{N} \norm{\vx}_4^4 - \frac{3}{N^2} \norm{\vx}_2^4.
\end{equation}
最后，对于随机向量，我们将其峰度定义为每个分量的标量峰度之和。
峰度是 ICA 的一个自然损失函数，因为对于高斯分布的 $X$，峰度为零；读者可以进一步验证伯努利-高斯分布具有正峰度。
因此，寻找非高斯独立分量的一个自然过程是搜索一组相互正交的方向 $\vV \in \R^{d \times k}$，使得
$\vV^\top \vX$ 具有最大峰度，其中 $\vX = \vU \vZ \in \R^{D \times N}$ 是伯努利-高斯
ICA 数据矩阵（使得 $\vV$ 在列排列的意义上近似 $\vU$）。
形式上，我们试图解决以下问题
\begin{equation}
    \max_{\vV^\top \vV = \vI} \kurt(\vV^\top \vX).
\end{equation}
在一个极端情况下，我们可以设置 $k = D$ 并试图一次性恢复整个字典 $\vU$。可以证明，这个问题可以用我们之前看到的 MSP 算法来解决。
在另一个极端情况下，我们可以设置 $k=1$ 并试图一次恢复一个单一方向（$\vU$
的列），在每一步之后、寻找另一个方向之前执行\textit{降维}（deflation），即用 $(\vI - \vu\vu^\top)
\vX$ 替换数据矩阵 $\vX$。
在 $k=1$ 增量方法的可扩展性与 $k=D$ 方法的效率和鲁棒性之间存在着自然的权衡。

\paragraph{增量 ICA：正确性与 FastICA 算法。}
由 Hyv\"{a}rinen 和 Oja \cite{hyvarinen-1997} 提出的 FastICA 算法，是一种用于求解
ICA 中 $k=1$ 峰度最大化方案的快速不动点算法。
手头的问题是
\begin{equation}\label{eq:kurtosis-maximization-sphere-finitesample}
    \max_{\norm{\vv}_2^2 = 1}\, \kurt(\vX^\top \vv).
\end{equation}
首先，我们将对该目标进行一些非常基础的分析以验证其正确性。请注意，通过变量替换 $\vw =
\vU^\top \vv$，该问题等价于
\begin{equation*}
    \max_{\norm{\vw}_2^2 = 1}\,
    \mathrm{kurt}(\vZ^\top \vw).
\end{equation*}
这个目标足够简单，以至于我们可以对其作为恢复字典 $\vU$ 的公式的正确性做出强有力的陈述。
例如，在 $N \to \infty$
的总体设定中，我们可以利用峰度的可加性（\Cref{exercise:kurtosis-linearity-properties}）以及我们对独立分量假设的归一化，将前一个问题等价地写为
\begin{equation}\label{eq:kurtosis-maximization-sphere-population-simple}
    \max_{\norm{\vw}_2^2 = 1}\,
    \sum_{i=1}^d \mathrm{kurt}(z_i) w_i^4.
\end{equation}
可以证明，在伯努利-高斯假设下，该问题的优化景观是“良好的”（\Cref{exercise:kurtosis-sphere-landscape}）——这意味着目标函数的所有局部极大值都对应于其中一个独立分量的恢复。
计算这些极大值之一的一种高效且可扩展的方法是通过一阶优化算法，该算法迭代地沿着目标函数的梯度方向，并投影到约束集 $\set{\vw
\given \norm{\vw}_2^2 = 1}$ 上。
由于我们假设每个 $z_i$ 满足 $\Var(z_i)=1$，对于较大的 $N$，我们有
\begin{equation}\label{eq:kurtosis-approximation-l4}
    \kurt(\vX^\top \vu)
    \approx
    \tfrac{1}{N} \norm{\vX^\top \vu}_4^4 - 3 \norm{\vu}_2^4.
\end{equation}
然后我们可以推导出梯度的相应近似值：
\begin{equation*}
    \nabla_{\vu} \kurt(\vX^\top \vu)
    \approx
    \tfrac{4}{N} \vX (\vX^\top \vu)^{\hada 3}
    - 12 \norm{\vu}_2^2 \vu.
\end{equation*}
FastICA 算法使用不动点方法来计算最大峰度的方向。它从峰度最大化问题的一阶最优性条件出发，给定前面的梯度近似和约束集，其形式为
\begin{align}\label{eq:kurtosis-max-sphere-stationarity}
    \vX (\vX^\top \vu)^{\hada 3}
    =
    \underbrace{
        \ip*{\vu}{
            \vX (\vX^\top \vu)^{\hada 3}
    }}_{\lambda} \vu,
\end{align}
其中 $\lambda$ 的具体值是使用 $\vu$ 上的单位范数约束确定的。
\Cref{exercise:sphere-calculus} 描述了从第一性原理推导这些最优性条件所需的数学背景。
方程 \eqref{eq:kurtosis-max-sphere-stationarity}
被峰度最大化问题的\textit{任何}临界点所满足；我们希望推导出一个仅被极大值点满足的方程。
在注意到 $\lambda = \norm{\vX^\top \vu}_4^4$ 之后，我们将
\eqref{eq:kurtosis-max-sphere-stationarity} 等价地重新表达为修改后的方程
\begin{align}\label{eq:kurtosis-max-sphere-stationarity-modified}
    \frac{1}{N}\vX (\vX^\top \vu)^{\hada 3}
    -
    3 \vu
    =
    \left(
        \frac{\lambda}{N} - 3
    \right)
    \vu,
\end{align}
并意识到 \eqref{eq:kurtosis-maximization-sphere-finitesample}
的任何极大值点都必须满足 $\lambda / N - 3 > 0$，假设 $N$ 足够大。
因此，我们可以对 \eqref{eq:kurtosis-max-sphere-stationarity-modified}
的两边进行\textit{归一化}，得到以下被
\eqref{eq:kurtosis-maximization-sphere-finitesample} 的每个极大值点所满足的不动点方程：
\begin{align}\label{eq:kurtosis-max-sphere-fxp}
    \frac{
        \frac{1}{N}\vX (\vX^\top \vu)^{\hada 3}
        -
        3 \vu
    }{
        \norm*{
            \frac{1}{N}\vX (\vX^\top \vu)^{\hada 3}
            -
            3 \vu
        }_2
    }
    =
    \vu.
\end{align}
迭代由该不动点表达式左侧定义的映射，即可得到 Hyv\"{a}rinen 和 Oja \cite{hyvarinen-1997} 的 FastICA 算法：
\begin{equation}
    \begin{split}\label{eq:fast-ica}
        \vv^+ &= \tfrac{1}{N}\vX (\vX^\top \vu)^{\hada 3}- 3 \vu
        ,  \\
        \vu^+ &= \vv^+ / \norm*{\vv^+}_2.
    \end{split}
\end{equation}
事实证明，FastICA 算法极快地（实际上以\textit{三次}收敛率）收敛到字典 $\vU$ 的列；感兴趣的读者可以查阅
\cite{hyvarinen-1997} 了解详情。
将 FastICA 算法 \eqref{eq:fast-ica} 与 \Cref{sub:pca} 中研究的用于 PCA 问题的幂法以及
MSP 算法 \eqref{eq:msp_iteration} 进行比较，我们注意到惊人的相似性。事实上，FastICA
本质上是一种修改后的幂法，它涉及经验峰度的梯度，而不是 PCA 目标中更简单的线性梯度。

\section{过完备低维子空间的混合}
\label{sec:dictionary_learning}
正如我们所见，完备字典学习享有优雅的计算理论，其中我们保持对称的自编码结构 $\cE(\vx) = \vU^\top
\vx$，$\cD(\vz) = \vU \vz$，并使用一种可扩展的类似幂法的算法（MSP 算法）从数据中学习正交字典/码本
$\vU$。然而，为了学习一般高维数据分布的表征，必须将码本的大小扩展到超出正交性要求的范围——这意味着我们必须有 $\vA \in \R^{D
\times m}$，其中 $m \gg
D$，对应于\textit{过完备}（overcomplete）字典/码本的情况，\footnote{我们在这里将符号从 $\vU$
更改为 $\vA$，以强调过完备字典 $\vA$ 的非正交性和非方形形状。} 以及信号模型
\begin{equation}\label{eq:model-DL-overcomplete}
    \vx =  \vA \vz + \veps,
    \quad
    \norm{\vz}_0 = d \ll m.
\end{equation}
转向过完备情况既有几何动机，也有物理/建模动机。
在几何上，回想一下，在我们最初从子空间混合数据模型到稀疏字典模型的简化中，$\R^D$ 中 $K$ 个子空间（每个维度为
$d$）的混合导致了形状为 $\vA \in \R^{D \times Kd}$ 的字典。
换句话说，过完备字典对应于\textit{更丰富}的子空间混合，具有更多不同的变异模式来对高维数据分布进行建模。
在建模方面，我们可以在真实世界数据上运行计算实验，以揭示过完备表征所赋予的额外建模能力。

\begin{figure}[t]
    \centering
    \begin{subfigure}{0.9\linewidth}
        \centering
        \includegraphics[width=\linewidth]{\toplevelprefix/chapters/chapter2/figs/msp_atoms_patches_new.png}
        \caption{}
    \end{subfigure}
    \begin{subfigure}{0.9\linewidth}
        \centering
        \includegraphics[width=\linewidth]{\toplevelprefix/chapters/chapter2/figs/palm_atoms_patches_new.png}
        \caption{}
    \end{subfigure}
    \caption{完备（正交）字典和过完备字典学习到的字典原子的比较，训练用于重建从 MNIST 数字中提取的 $8 \times
        8$ 图像块。两个字典都在 $10^4$ 个具有非平凡内容的随机图像块上训练了 $6000$ 个 epoch，并且使用 LASSO
        目标和 $\lambda=0.1$ 计算稀疏代码（见
        \eqref{eq:sparse_dl_lasso}）。颜色映射中黑色表示负值，白色表示正值。\textbf{上图：}使用 MSP
        算法 \eqref{eq:msp_iteration} 学习的正交字典被限制为不超过 $64$
        个原子；学习到的原子大致对应于“尖峰和平板”（spike and
        slab）字典，并且在保留的测试数据上实现了相对较差的重建稀疏度水平（相对于 $10^{-1}$ 的阈值，代码平均约为
        $17$-稀疏）。\textbf{下图：}相比之下，过完备字典（此处有 $8^3$ 个原子；我们可视化了 $64$
        个随机子集）学习到了具有语义意义的字典原子，对应于带符号的定向边缘，这些边缘可以拼凑在一起创建数字图像块，并实现卓越的重建和稀疏度水平。代码平均约为
        $20$-稀疏，同时比正交字典的代码大 $8$ 倍。为了计算字典，我们在 \eqref{eqn:DL-overcomplete}
    的适当正则化版本上使用了基于近端交替线性化最小化（proximal alternating linearized minimization）的优化器。}
    \label{fig:ReconMNIST}
\end{figure}

\begin{example}
    给定手写数字的采样图像，\Cref{fig:ReconMNIST}(a) 显示了将正交字典拟合到数据集的结果。
    相比之下，\Cref{fig:ReconMNIST}(b) 显示了在这些样本上运行用于学习过完备字典的优化算法（我们将在本章后面详细介绍）的结果。
    请注意，表征变得更加稀疏，并且码本变得更具可解释性——它们由构成数字笔画的基本图元（即定向边缘）组成。
\end{example}

事实上，过完备字典学习最初是作为一种生物学上合理的图像表征算法被提出的，其基础是关于视觉皮层早期阶段如何表征视觉刺激的经验证据
\cite{Olshausen1996-ap,Olshausen1997-yv}。

在本节的其余部分，我们将概述过完备字典学习的概念和计算基础。
假设模型 \eqref{eq:model-DL-overcomplete} 满足稀疏代码 \(\vz\)、过完备字典 \(\vA\)
和稀疏度水平 \(d\)，并且给定 \(\vx\) 的样本 \(\vX = [\vx_1, \dots,
\vx_N]\)，我们希望学习一个编码器 \(\cE : \R^D \to \R^m\) 将每个 \(\vx\) 映射到其{稀疏代码}
\(\vz\)，以及一个解码器 \(\cD(\vz) = \vA \vz\) 从其稀疏代码重建每个 \(\vx\)。 %
用图表表示为：
\begin{equation}
    \x \xrightarrow{\hspace{4mm} \mathcal{E} \hspace{4mm}}  \z
    \xrightarrow{\hspace{2mm} \mathcal{D} = \vA \hspace{2mm}}   \hat{\x}.
    \label{eqn:autoencoding-DL-overcomplete}
\end{equation}

我们将从编码器 $\cE$ 的构建开始。
我们将逐步进行：首先，\textit{给定真实字典
$\vA$}，我们将展示\textit{稀疏编码}问题如何提供一种优雅、可扩展且可证明正确的算法，用于恢复 $\vx$ 的稀疏代码 $\vz$。
尽管这个问题在最坏情况下是 NP 难的，但对于\textit{非相干}（incoherent）的字典
$\vA$（即列之间不太相关的字典），它可以被高效且可扩展地解决。
该解决方案所包含的编码器架构将不再是对称的：我们将看到它具有原始深度网络的形式，该网络依赖于字典 $\vA$。

然后我们将继续学习解码器 $\cD$ 的任务，或者等价地学习过完备字典 $\vA$。
我们将推导一种用于过完备字典学习的算法，该算法允许我们利用稀疏编码的思想同时学习码本 $\vA$ 和稀疏代码 $\vz$。
最后，我们将讨论关于可学习稀疏编码的更现代的视角，这将引导我们走向完全不对称的编码器-解码器结构，作为
\eqref{eqn:autoencoding-DL-overcomplete} 的替代方案。
在这里，解码器将对应于稀疏字典学习问题的增量解决方案，并产生一对用于稀疏字典学习的类似深度网络的编码器-解码器。
这种结构将预示着本专著其余部分中即将出现的许多发展，因为我们将从分析模型过渡到现代神经网络。

\subsection{使用过完备字典的稀疏编码}

在本节中，我们将考虑数据模型
\eqref{eq:model-DL-overcomplete}，它容纳了许多基序（或\textit{原子}）的稀疏线性组合。给定满足该模型的数据
\(\{\vx_{i}\}_{i =
    1}^{N} \subseteq
\R^{D}\)，即可以表示为
\begin{equation}\label{eq:vectorized_sparse_dl_dgp}
    \vx_{i} = \vA\vz_{i} + \veps_{i}, \qquad \forall i \in [N]
\end{equation}
对于具有 $m$ 个原子的某个字典 $\vA \in \bR^{D \times m}$、满足 $\norm{\vz_i}_0 \leq
d$ 的稀疏代码 $\vz_i$ 以及小误差 $\veps_i$，稀疏编码问题是在给定字典 $\vA$ 的情况下，尽可能准确地从数据
$\vx_i$ 中恢复代码 $\vz_i$。
当字典 \(\vA\) 是\textit{非相干}的（即内积 \(\va_{i}^{\top}\va_{j}\)
一致较小，因此原子几乎正交）时，解决该问题的高效算法能够成功。\footnote{事实证明，在高维空间中，很容易打包数量远大于环境维度的近乎正交的向量
\cite{Wright-Ma-2022}。 }

请注意，我们可以将 \(\vx_{i}\) 收集为 \(\vX = \mat{\vx_{1},
\dots, \vx_{N}} \in \R^{D \times N}\)，将 \(\vz_{i}\) 收集为 \(\vZ =
\mat{\vz_{1}, \dots, \vz_{N}} \in \R^{d  \times N}\)，并将 \(\veps_{i}\)
收集为 \(\vE = \mat{\veps_{1}, \dots,
\veps_{N}} \in \R^{D \times N}\)，从而将 \eqref{eq:vectorized_sparse_dl_dgp} 重写为
\begin{equation}\label{eq:sparse_dl_dgp}
    \vX = \vA\vZ + \vE.
\end{equation}
解决稀疏编码问题的一种自然方法是寻找与我们的观测结果一致的最稀疏信号，这自然引出了以下优化问题：
\begin{equation}\label{eq:sparse_dl_lasso}
    \min_{\vZ \in \R^{d \times N}}\bc{\norm{\vX - \vA\vZ}_{F}^{2} +
    \lambda \norm{\vZ}_{1}},
\end{equation}
其中逐元素计算的 \(\ell^1\) 范数 \(\norm{\vZ}_{1}\) 已知可以促进解的稀疏性 \cite{Wright-Ma-2022}。
这个优化问题被称为 LASSO 问题。

然而，与 PCA 或完备字典学习问题不同，没有明确的幂迭代类型算法来恢复
\(\vZ^{\star}\)。一个自然的替代方案是考虑使用梯度下降类型的算法来解决上述优化问题。
令 \(f(\vZ) = \norm{\vX - \vA\vZ}_{2}^{2} + \lambda \norm{\vZ}_{1}\)。
从概念上讲，我们可以尝试通过以下迭代来找到 \(\vZ^{\star}\)：
\begin{equation}
    \vZ_{t+1} \leftarrow \vZ_t + \eta \nabla f(\vZ_t).
\end{equation}
然而，由于与 \(\ell^1\) 范数 \(\norm{\vZ}_{1}\)
相关的项是非平滑的，我们不能直接运行梯度下降。对于这种类型的函数，我们需要用一种推广了梯度概念的东西（称为次梯度 \(\partial
f(\vZ)\)）来替换梯度 \(\nabla f(\vZ)\)：
\begin{equation}
    \vZ_{t+1} \leftarrow \vZ_t + \eta \partial f(\vZ_t).
\end{equation}
然而，众所周知，次梯度下降的收敛速度通常非常慢。因此，对于这种类型的优化问题，通常采用所谓的 {\em 近端梯度下降}（proximal
gradient descent）类型的算法。我们在 \Cref{subsec:pgd} 中对该方法进行了技术概述。

将近端梯度应用于 LASSO 目标函数 \eqref{eq:sparse_dl_lasso}
会引出经典的\textit{迭代收缩阈值算法}（iterative shrinkage-thresholding
algorithm，ISTA），它实现了以下迭代
\begin{eqnarray}
    \vZ_{1} &\sim& \dNorm(\vzero, \vI), \\
    \vZ_{t + 1} &=& S_{\eta\lambda}\rp{\vZ_{t} - 2\eta
    \vA^{\top}(\vA\vZ_{t} - \vX)}, \quad \forall t \geq 1,\label{eq:ista_update}
\end{eqnarray}
其中步长 \(\eta \geq 0\)，软阈值算子 \(S_{\alpha}\) 在标量上定义为
\begin{align}
    S_{\alpha}(x) &
    \doteq
    \begin{cases}x - \alpha, & x \geq \alpha, \\ 0,
        & -\alpha < x < \alpha, \\ x + \alpha, & x \leq -\alpha
    \end{cases}
    \\
    &=
    \sign(x) \max \set{\abs{x} - \alpha, 0},
\end{align}
并逐元素应用于输入矩阵。作为应用于凸问题的近端梯度下降算法，它保证收敛到全局最优解，并且可以直接推导出精确的收敛率 \cite{Wright-Ma-2022}。

\begin{figure}
    \centering
    \begin{tikzpicture}
        \begin{axis}[
                width=0.6\linewidth,
                axis lines=middle,
                xmin=-3, xmax=3,
                ymin=-2, ymax=2,
                xtick=\empty,
                ytick=\empty,
            ]
            \addplot[thick, domain=-3:-1] {x + 1};
            \addplot[thick, domain=-1:1] {0};
            \addplot[thick, domain=1:3]  {x - 1};

            \addplot[only marks, mark=*, mark size=1.5pt]
            coordinates {(-1,0) (1,0)};

            \node at (axis cs:1,-0.35) {$\lambda$};
            \node at (axis cs:-1,-0.35) {$-\lambda$};

        \end{axis}
    \end{tikzpicture}
    \caption{\small 软阈值函数 \(S_{\lambda}\) 的图像。在 \(-\lambda\) 和
    \(+\lambda\) 之间，函数值被截断为 \(0\)，而较大输入的（绝对）值则简单地收缩 \(\lambda\)。}
\end{figure}

我们现在花点时间来评论一下 \eqref{eq:ista_update} 中更新算子的函数形式。它采取的形式为
\begin{equation}
    \vZ_{t + 1} = \texttt{nonlinearity}(\vZ_{t} +
    \texttt{linear}^{\top}(\texttt{linear}(\vZ_{t}) + \texttt{bias})).
\end{equation}
这种函数形式与残差网络层的形式非常相似，即
\begin{equation}
    \vZ_{t + 1} = \vZ_{t} +
    \texttt{linear}_{1}^{\top}(\texttt{nonlinearity}(\texttt{linear}_{2}(\vZ_{t})
        + \texttt{bias}_{1}) + \texttt{bias}_{2},
    \end{equation}
    只是解耦了线性映射（即矩阵乘法），添加了偏置，并移动了非线性。
    \eqref{eq:ista_update} 中的非线性与深度学习中常用的 ReLU 激活函数 $x \mapsto
    \max\set{x, 0}$ 非常相似——特别是，软阈值算子就像一个“带符号的” ReLU 激活，它也被偏置平移了。
    因此，ISTA 可以被视为一个原始（循环单层）神经网络的前向传递。我们在 \Cref{ch:deep-networks}
    中论证了此类操作对于深度表征学习是必不可少的。

    \subsection{过完备字典学习}

    回想一下，我们有数据模型
    \begin{equation}\label{eq:overcomplete-dl-section-data-model}
        \vX = \vA\vZ + \vE,
    \end{equation}
    其中 \(\vZ\) 是稀疏的，我们之前的目标是在已知数据 \(\vX\) 和字典原子 \(\vA\) 的情况下估计 \(\vZ\)。
    现在我们转向更实际也更困难的情况，即我们既不知道 \(\vA\) 也不知道 \(\vZ\)，并试图从大型数据集中学习它们。

    对 \Cref{eq:sparse_dl_lasso} 的直接推广建议解决以下问题
    \begin{equation}
        \min_{\tilde{\vA}, \tilde{\vZ}}\bc{\norm{\vX -
        \tilde{\vA}\tilde{\vZ}}_{F}^{2} + \lambda \norm{\tilde{\vZ}}_{1}}.
        \label{eqn:DL-overcomplete-try}
    \end{equation}
    然而，\Cref{eqn:DL-overcomplete-try} 中的双线性项引入了阻碍收敛的尺度模糊性：给定任意点 $(\tilde{\vA},
    \tilde{\vZ})$ 和任意常数 $c>0$，代换 $(c^{-1}\tilde{\vA}, c\tilde{\vZ})$ 给出的损失值为
    \begin{equation}
        \norm{\vX - \tilde{\vA}\tilde{\vZ}}_{F}^{2} + c\lambda
        \norm{\tilde{\vZ}}_{1}.
    \end{equation}
    显然，通过取 $c \to 0$，该损失在 $c$ 上被最小化，这对应于使重新缩放的字典 $c^{-1} \tilde{\vA}$
    趋于“无穷大”。特别是，求解 \eqref{eqn:DL-overcomplete-try} 的任何优化算法的迭代都不会收敛。

    \eqref{eqn:DL-overcomplete-try} 的这个问题可以通过对字典 $\tilde{\vA}$ 的列的尺度添加约束来解决。
    例如，通常假设 \eqref{eq:overcomplete-dl-section-data-model} 中字典 $\vA$
    的每一列 $\vA_j$ 具有有界的 $\ell^2$ 范数——例如，不失一般性地，界限为 $1$。
    然后我们将其作为约束强制执行，得到目标
    \begin{equation}
        \min_{\tilde{\vZ}, \tilde{\vA} \,:\, \norm{\tilde{\vA}_j}_2 \leq 1}
        \bc{\norm{\vX - \tilde{\vA}\tilde{\vZ}}_{F}^{2} + \lambda
        \norm{\tilde{\vZ}}_{1}}.
        \label{eqn:DL-overcomplete}
    \end{equation}
    诸如 \eqref{eqn:DL-overcomplete} 之类的约束优化问题可以通过许多复杂的算法来解决
    \cite{nocedal2006numerical}。然而，我们上一节中用于解决稀疏编码问题的同一个近端梯度下降算法实际上提供了一个简单且可扩展的算法。
    我们可以通过约束集的特征函数将每个约束编码为额外的正则化项——详细信息在
    \Cref{example:prox-of-characteristic-function} 中给出。
    将近端梯度下降应用于由此产生的正则化问题等价于\textit{投影梯度下降}（projected gradient
    descent），其中，在每次迭代中，采取梯度下降步骤后的迭代值被投影到约束集上。

    \begin{remark}[$\ell^4$ 最大化与 $\ell^1$ 最小化]
        请注意，上述问题公式自然地源于稀疏编码的 LASSO 公式 \eqref{eq:sparse_dl_lasso}。我们通过
        \(\ell^1\) 范数来促进解的稀疏性。
        然而，如果我们只对恢复过完备字典 \(\vA\) 感兴趣，\Cref{sec:complete-dictionary}
        中介绍的 \(\ell^4\) 最大化方案也可以推广到过完备情况，而无需任何重大修改。感兴趣的读者可以参考
        \cite{Qu2020Geometric} 的工作。
    \end{remark}

    上述问题 \eqref{eqn:DL-overcomplete}（我们称之为\textit{过完备字典学习}）是非凸的，因为这里的
    \(\vA\) 和 \(\vZ\)
    都是未知的。它不能轻易地用标准的凸优化工具包来解决。尽管如此，由于它很有趣、陈述简单且在实践中很重要，因此有许多重要的工作致力于该问题的不同算法和理论分析。
    在这里，出于本手稿的兴趣，我们提出了一种更接近深度学习精神的惯用方法来解决这个问题。

    根据我们上面处理 LASSO 问题的经验，很容易看出，对于两个未知数 \(\vA\) 和
    \(\vZ\)，如果我们固定一个并优化另一个，每个子问题实际上都是凸的且易于求解。这自然表明我们可以尝试通过交替最小化
    \(\vZ\) 或 \(\vA\)（例如使用梯度下降）来求解上述程序 \eqref{eqn:DL-overcomplete}。
    结合自然的初始化选择，这导致了以下迭代方案：
    \begin{align}
        &\vZ^{\ell + 1} = S_{\eta\lambda}\rp{\vZ^{\ell}
        - 2\eta\vA_{+}^{\top}(\vA_{+}\vZ^{\ell} - \vX)},
        \quad \vZ^1 = \Zero,
        \quad \forall \ell \in [L] \label{eqn:ISTA-update}\\
        &\vZ^+ = \vZ^L,\\
        &\vA_{t + 1} = \proj_{\norm{(\,\cdot\,)_j}_2 \leq 1,\,\forall j}\left(
            \vA_{t} - 2\nu (\vA_{t}\vZ^+ - \vX)(\vZ^+)^\top
        \right), \nonumber \\
        &\quad\mbox{其中} \quad (\vA_{1})_j \simiid \dNorm(\vzero,
        \tfrac{1}{D}\vI), \enspace \forall
        j \in [m],
        \quad\;\; \forall t \in [T],\label{eqn:DL-update}\\
        &\vA_+ = \vA_T,
    \end{align}
    其中 $\vA$ 更新中的投影操作通过 $\vA_j \mapsto \vA_j / \max
    \set{\norm{\vA_j}_2,
    1}$ 确保每列最多具有单位 $\ell^2$ 范数，并且其中 $\vA_+$ 的初始化是每列独立同分布于 $\cN(\Zero,
    \tfrac{1}{D} \vI)$。
    上述内容构成了一个交替最小化的“块”，我们重复执行这样的块，每个块都有独立的初始化，直到收敛。
    在上面，我们使用了两个独立的索引 $\{t\}$ 和 $\{\ell\}$
    来指示迭代。正如我们稍后将看到的，这允许我们在深度学习的背景下分别解释这两个更新。

    尽管字典学习问题是一个非凸问题，但已经证明交替最小化类型的算法确实能收敛到正确的解，至少在局部是这样。例如，参见 \cite{alekh-2016}。
    作为实际演示，上述算法（其中 $L = T =
    1$）被用于生成 \Cref{fig:ReconMNIST} 中过完备字典学习的结果。

    \subsection{学习的深度稀疏编码}
    \label{sec:LISTA}
    上一节中用于过完备字典学习的交替最小化算法（\Cref{eqn:ISTA-update,eqn:DL-update}）的主要见解是注意到\textit{当我们固定
        \(\vA\) 时，$\vZ^\ell$ 的 ISTA 更新 \eqref{eqn:ISTA-update}
        看起来就像是一个权重由 \(\vA\)（和
    \(\vA^{\top}\)）给出的深度神经网络的前向传递}。但在通常情况下，我们不知道真实的 $\vA$，并且当前的估计
    $\vA_+$ 可能是错误的。因此，需要基于使用当前稀疏代码估计 $\vZ^+$ 重建 $\vX$ 的残差，使用
    \eqref{eqn:DL-update}
    对其进行进一步更新。交替最小化算法迭代这两个过程直到收敛。但我们可以进行外推，通过将这些见解与深度学习技术相结合来设计其他学习过程。这导致了更具可解释性的网络架构，这将是贯穿本手稿的一个反复出现的主题。

    \paragraph{学习的 ISTA（Learned ISTA）。}
    上述对交替最小化的深度网络解释更多是概念性的而非实用性的，因为该过程可能相当低效，并且需要许多层或迭代才能收敛。
    但这主要是因为我们试图从 \(\vX\) 中同时推断 \(\vZ\) 和 \(\vA\)。
    在\textit{监督}设定下，这个问题可以被显著简化，并且上述迭代可以变得更加高效，在这种设定下，我们有一个根据
    \eqref{eq:overcomplete-dl-section-data-model} 分布的输入和输出对 \((\vX,
    \vZ)\) 的数据集，并且我们只试图为逐层可学习的稀疏编码迭代 \eqref{eq:ISTA-update-layerwise}
    学习 \(\vA^\ell\)：
    \begin{align}
        & \vZ^{\ell + 1} = S_{\eta\lambda}\rp{\vZ^{\ell} -
        2\eta(\vA^{\ell})^{\top}(\vA^\ell\vZ^{\ell} - \vX)}, \quad
        \forall \ell \in [L].
    \end{align}
    如果我们将每次迭代的算子表示为 $\vZ^{\ell + 1} =
    f(\vA^\ell, \vZ^\ell)$，上述迭代可以用图表说明：
    \begin{equation*}
        \vX, \vZ^1 \xrightarrow{\hspace{1mm} f(\vA^1,\,\cdot\,)
        \hspace{1mm}}  \vZ^2 \xrightarrow{\hspace{1mm}
        f(\vA^2,\,\cdot\,) \hspace{1mm}}  \vZ^3
        \xrightarrow{\hspace{1mm} f(\vA^3,\,\cdot\,) \hspace{1mm}}
        \cdots \vZ^{L}  \xrightarrow{\hspace{1mm}
        f(\vA^{L},\,\cdot\,) \hspace{1mm}} \vZ^{L + 1} \approx \vZ.
        \label{eqn:deep-sparse-encoding}
    \end{equation*}
    因此，给定这种序列架构，为了学习每一层的算子 \(\vA^\ell\)，通过最小化最终代码 $\vZ^L$ 与真实值 $\vZ$
    之间的误差来学习它（例如通过反向传播（BP）\footnote{有关 BP 的简要描述，请参见
    \Cref{sec:autodiff}。}）是完全自然的：
    \begin{equation}\label{eq:lista-training}
        \min_{\{\vA^\ell\}} \big\|\vZ^L(\vA^1, \ldots, \vA^{L}) - \vZ\big\|_2^2.
    \end{equation}
    这是学习的 ISTA（Learned ISTA，LISTA）算法 \cite{gregor2010learning}
    的基础，它可以被视为深度神经网络的学习算法，该网络试图模拟从 $\vX$ 到 $\vZ$
    的稀疏编码过程。特别是，它可以被视为一种\textit{简单的表征学习算法}。事实上，同样的方法论可以作为理解在更强大的网络架构（如
    Transformer）中计算的表征的基础。我们在 \Cref{ch:deep-networks} 中详细展开了这些含义。

    \paragraph{稀疏自编码器（Sparse Autoencoders）。}
    过完备字典学习的最初动机是为高维数据提供一个简单的生成模型。通过 LISTA
    我们已经看到，此外，用于学习稀疏使用的过完备字典的迭代算法为类似 ReLU 的深度网络提供了解释，我们将在后面的章节中将其推广到比
    \eqref{eq:overcomplete-dl-section-data-model} 更复杂的数据分布。
    但同样值得注意的是，即使在大模型的现代时代，数据生成模型
    \eqref{eq:overcomplete-dl-section-data-model}
    也为\textit{解释预训练的大规模深度网络（如
    Transformer）中的特征}提供了有用的实践基础，其遵循的假设是，这些网络中的（\textit{先验}不可解释的）特征由底层特征的稀疏“叠加”（superpositions）组成，而这些底层特征本身是可解释的
    \cite{elhage2022superposition}。这些\textit{无监督}学习范式通常也比 LISTA
    对数据更友好，因为 LISTA 需要大量标记的 $(\vX, \vZ)$ 对进行监督训练。

    我们可以利用上面对 LISTA 算法的推导来理解该研究领域的常见实践。
    在最直接的实例化中（参见 \citep{huben2024sparse,
    gao2025scaling}），从预训练深度网络 $h$ 中收集大量来自不同输入 $\vx_i$
    的特征，这些输入本身是根据所需的解释任务选择的。\footnote{例如，输入 $\vx_i$
        可以对应于包含不同编程语言的计算机代码样本的文本，我们的任务是试图在 Transformer 特征图 $h$
        中识别出对应于输入的不同显著方面的可解释特征，例如特定的编程语言（在不同的输入“类别”中有所不同）或在当前位置插入匹配括号的需要（在不同的输入“类别”中很常见）。我们在
    \Cref{ch:applications} 中更详细地讨论了深度网络（特别是 Transformer）在文本表征学习中的使用。}
    为简单起见，我们将使用 $h$ 来表示所讨论的预选特征图，具有 $D$ 维特征；给定 $N$ 个样本输入，令 $\vH \in \bR^{D
    \times N}$ 表示 $h$ 的完整特征矩阵。
    然后，通过 LASSO 损失 \eqref{eq:sparse_dl_lasso} 训练一个所谓的稀疏自编码器 $f :
    \bR^D \to \bR^d$，其解码器为 $g : \bR^d \to
    \bR^D$：
    \begin{equation}\label{eq:sae-loss}
    \min_{f, g} \| \vH - g(f(\vH))) \|_F^2 + \lambda \|f(\vH)\|_1 ,
\end{equation}
其中稀疏自编码器 $f$ 采用单层神经网络的形式，即 $f(\vh_i) = \sigma(\vW_{\mathrm{enc}}(\vh_i - \vb)
+ \vb_{\mathrm{enc}})$，其中 $\sigma(x) = \max \{x, 0\}$ 是 ReLU
激活函数，而解码器 $g$ 是线性的，因此 $g(\vz_i)
= \vW_{\mathrm{dec}}\vz + \vb$。\footnote{由于 \(f\)
    是一个稀疏化算子，我们将其表示为类似于之前使用的 \(f\)，例如在 LISTA 中，\(f\)
    表示对输入进行稀疏编码的优化算法。在整本书中，\(f\)
将被用来表示这种稀疏的以及其他结构化的变换，包括\textit{深度网络编码器}的变换（\Cref{ch:lossy-compression,ch:deep-networks}）。}

参数化和训练过程 \eqref{eq:sae-loss}
最初可能看起来像是深度学习在稀疏编码问题上的任意应用，但它实际上与我们上面研究的具有学习字典的逐层稀疏编码算法高度一致。
特别是，回想一下 LISTA %
架构 $\vZ^{L} = f(\vA^L, f(\vA^{L-1}, \cdots, f(\vA^1, \vX) \cdots ))$。
在 $L=2$ 的特殊情况下，我们有
\begin{equation}
    \vZ^2 = f(\vA^1, \vX)
    =S_{\eta\lambda}\rp{\vZ^{1} - 2\eta(\vA^{1})^{\top}(\vA^1\vZ^{1} - \vX)}.
\end{equation}
让我们假设所讨论的稀疏代码 $\vZ$ 是非负的，即 $\vZ \geq \Zero$。\footnote{在数据生成模型
    \eqref{eq:overcomplete-dl-section-data-model} 中，只需将 $\vA$
    中的列数加倍，就可以将任意字典和稀疏代码对 $(\vA, \vZ)$ 替换为满足 $\vZ \geq \Zero$
的对，因此从建模的角度来看，这是一个非常温和的假设。}
那么（参见 \Cref{example:prox-of-nonnegative-l1}），我们可以考虑从带有对 $\vZ$
的额外非负性约束的稀疏编码目标中获得的等效 LISTA 架构，即
\begin{equation}
    \vZ^2 = f(\vA^1, \vX)
    =\max \set*{\vZ^{1} - 2\eta(\vA^{1})^{\top}(\vA^1\vZ^{1} - \vX)
    - \lambda\eta \mathbf{1}, 0},
\end{equation}
经过一些代数运算后，将其表示为
\begin{equation}
    \vZ^2 = f(\vA^1, \vX)
    =\max \set*{
        2\eta(\vA^{1})^{\top}
        +
        \left(
            \vZ^{1} - 2\eta(\vA^{1})^{\top}\vA^1\vZ^{1} - \lambda\eta \mathbf{1}
        \right), 0
    }.
\end{equation}
鉴于能够将稀疏代码初始化 $\vZ^1$ 作为可学习参数进行更改（在当前框架中，其所有列必须等于同一个可学习向量），这具有带可学习偏置的
ReLU 神经网络的形式——与稀疏自编码器 $f$ 完全相同！
此外，为了\textit{解码}学习到的稀疏代码 $\vZ^2$，自然会应用学习到的字典 $\vZ^2 \mapsto \vA^1
\vZ^2$。那么这与 SAE 解码器 $g$ 之间的唯一区别就是额外的偏置 $\vb$，在技术上，它可以被吸收到训练目标
\eqref{eq:sae-loss} 中的 $\vH$ 和 $f$ 中。

因此，SAE 的参数化和训练过程与 $L=1$ 的 LISTA 训练以及修改后的训练目标相吻合——使用 LASSO 目标
\eqref{eq:sparse_dl_lasso}，它仍然 {\em
无监督的}，而不是原始LISTA中使用的有监督重建损失\eqref{eq:lista-training}。特别是，我们可以根据我们在\eqref{eq:ISTA-update-layerwise}中对LISTA架构的逐层稀疏编码的解释来理解SAE架构。这种联系暗示了许多用于改进实际可解释性方法的新设计策略，其中许多策略仍然诱人地未被探索。我们在\Cref{ch:consistent-representations}中开始展示与更广泛的自编码方法的一些联系。

\paragraph{逐层学习的稀疏编码？}

在有监督的设置中，LISTA提供了稀疏编码迭代的深度神经网络模拟，具有逐层学习的字典，其灵感来自交替最小化；即使在无监督的设置中，同样的方法也可以应用于学习，就像稀疏自编码器一样。但是，寻求低维结构的优化算法与深度网络架构之间的联系远不止于此，它暗示了一系列可扩展且自然的神经学习架构，这些架构甚至可能在不使用反向传播的情况下使用。

作为一个简单的说明，我们回到交替最小化迭代\eqref{eqn:ISTA-update}和\eqref{eqn:DL-update}。该方案在每次此类更新时随机重新初始化字典$\vA_1$。一种改进方法是使用\textit{热启动}（warm
starting），即使用先前对字典的估计$\vA_+$来生成残差。如果我们随后将每次ISTA更新\eqref{eqn:ISTA-update}视为一个层，并允许相关的字典（现在与稀疏代码更新耦合为$\vA_\ell$）随时间更新，这将导致一种“逐层可学习”的稀疏编码方案：
\begin{align}
    & \vZ^{1}
    = \Zero, \quad (\vA_{1})_j
    \simiid \dNorm(\vzero, \tfrac{1}{D}\vI), \enspace \forall j \in [m], \\
    & \vZ^{\ell + 1} = S_{\eta\lambda}\rp{\vZ^{\ell}
        - 2\eta(\vA_{\ell})^\top(\vA_{\ell}\vZ^{\ell}
    - \vX)},\label{eq:ISTA-update-layerwise}\\
    & \vA_{\ell + 1} = \vA_{\ell}- 2\nu (\vA_{\ell}\vZ^{\ell + 1}
    - \vX)(\vZ^{\ell + 1})^\top.
\end{align}
请注意，此迭代 %
对应于在无限多个块上对$T = L = 1$的\eqref{eqn:ISTA-update}和\eqref{eqn:DL-update}的重新标记。
更新$\vZ$的每个“内部”步骤可以被视为原始深度神经网络的单层前向传播，而更新$\vA$的每个“外部”步骤可以被视为单层反向传播。特别是，该算法是前向优化和反向学习之间表现出明确分界的最简单情况。这种对数据$\vX$的短暂“特征”$\vZ$的前向更新与对持久参数$\vA$的反向更新之间的区别，在当前的神经网络和自编码器中仍然可以观察到——我们将在\Cref{ch:deep-networks}和\Cref{ch:consistent-representations}中对此进行更多讨论。

请注意，上述逐层方案还提出了一种合理的替代方案，以取代当前主要依赖于\Cref{app:BP-section}中详述的反向传播（BP）的端到端优化策略。将大型网络的训练从BP中解放出来将是未来最大的挑战和机遇之一，正如我们将在本书末尾的\Cref{ch:future}中进一步讨论的那样。

\section{总结与注释}

\paragraph{关键信息。}
我们在本章中介绍的理想化模型——PCA、ICA和字典学习——是在二十世纪发展起来的。关于每种方法都有许多专著，因此我们在这里只试图对关键概念和结果进行广泛的概述。正如我们将在接下来的章节中看到的，尽管这些可以说是用于低维数据分布的相当理想化的简单模型，但识别此类分布的基本思想甚至由此产生的算法都可以自然地推广或扩展到任意低维分布——这也是本书的主题。此外，正如我们将在接下来的章节中看到的，这些基本（线性）模型在很大程度上也可以作为局部原型，以有效地近似建模具有非线性低维结构的一般分布。

\paragraph{更多历史与参考文献。}
Jolliffe \cite{Jolliffe1986} 将主成分分析归功于 Pearson
\cite{Pearson1901}，以及独立提出该方法的 Hotelling
\cite{Hotelling1933}。在数学中，关于酉不变范数下低秩近似相关问题的主要结果归功于 Eckart 和 Young
\cite{Eckart1936-ep}，而其完全一般性的结果则归功于 Mirsky \cite{Mirsky1960-ek}。PCA
作为无监督表征学习中最简单的模型问题，继续在研究中发挥着重要作用：早在20世纪80年代，\textcite{Oja1982SimplifiedNM}
和 \textcite{Baldi89}
等工作就利用该问题来理解原始神经网络中的学习，而最近，它又被用作理解更复杂的表征学习框架（如扩散模型 \cite{wang2024diffusion}）的工具。

独立成分分析由 \textcite{Ans-1985} 提出，并由 Aapo Hyv\"{a}rinen
在20世纪90年代和21世纪初的一系列有影响力的著作中开创：总结见
\textcite{Hyvrinen-2000}。作为实际数据中出现的一种简单结构模型，它最初在盲源分离等应用中得到了广泛使用，其中每个独立成分
$z_i$ 代表一个独立源（例如音乐录音中与不同乐器相关的声音），这些源叠加产生观测值 $\vx = \vU \vz$。

在完备或正交的情况下，字典学习问题可以被视为二十世纪信号处理的基础问题之一，特别是在线性系统理论中，傅里叶基在其中起着关键作用；从20世纪80年代开始，计算调和分析领域围绕着研究此类替代字典而形成，用于那些只能在非傅里叶基（例如小波）中实现最佳近似的信号类
\cite{Donoho1998-zf}。然而，冗余基或过完备字典情况的重要性，直到 Olshausen 和 Field 的开创性工作
\cite{Olshausen1996-ap,Olshausen1997-yv}
之后才被凸显出来。早期的后续工作为学习稀疏使用的过完备字典奠定了概念和算法基础，通常旨在表示自然图像
\cite{Donoho2001-tl,DonohoD2003-PNAS,Elad2006-yi,Murray2007-cw,aharon2006k,Mairal2014-nq,Gribonval2015-fy}。后来，作为无监督表征学习的一个重要且非平凡的模型问题，该问题引起了大量的理论兴趣，导致信号处理、理论机器学习和理论计算机科学社区对其进行研究，特别关注在何种条件下该问题可以被可证明且有效地解决。这一系列中值得注意的工作（不详尽）包括
\textcite{Spielman2012-tl} 和 \textcite{sun2017complete_a}
关于完备情况的工作；\textcite{pmlr-v40-Arora15}；\textcite{Barak2015-uu}；以及
\textcite{Qu2020Geometric}。关于这个陈述简单的问题，许多深刻的理论问题仍然悬而未决，部分原因可能是由于该问题在最坏情况下的
NP 困难性（例如，参见 \textcite{Tillmann2015-up}）。

\begin{table}[tb]
    \centering
    \caption{本章介绍的（广义）幂法总结}
    \label{tab:power-method-summary}
    \resizebox{\textwidth}{!}{
        \begin{tabular}{llll}
            \toprule
            问题 & 算法 & 迭代 & 强制执行的结构类型 \\
            \midrule
            PCA & 幂法 & $\vu_{t} = \frac{\vX\vX^\top
            \vu_{t}}{\norm{\vX\vX^\top \vu_{t}}_2}$ &
            $1$维子空间（单位向量） \\
            \addlinespace
            ICA & FastICA &
            $\vu_{t+1} = \frac{\tfrac{1}{N}\vX (\vX^\top \vu_{t})^{\hada 3}-
            3 \vu_{t}}{\norm*{\tfrac{1}{N}\vX (\vX^\top \vu_{t})^{\hada 3}-
            3 \vu_{t}}_2}$
            & $1$维子空间（单位向量） \\
            \addlinespace
            完备字典学习 & MSP算法 &
            $\vU_{t+1} = \mathcal{P}_{\mathrm O(D)}[( {\vU_t
            \vX})^{\hada 3}\vX^\top]$
            & $D$维子空间（正交矩阵） \\
            \bottomrule
        \end{tabular}
    }
\end{table}

从对这些用于低维结构的经典分析模型的研究中，我们希望强调的一点是各种\textit{广义幂法}（generalized power
methods）所发挥的共同作用——这些算法至少在局部能够非常快速地收敛到各种类型的低维结构。这类算法的术语遵循了
\textcite{JourneeM2010} 的工作。在较高层面上，模仿用于计算半正定矩阵 $\vA \in \bR^{n \times
n}$ 的最大特征向量的经典幂迭代，即
\begin{equation}
    \vu_{t+1} = \frac{\vA \vu_{t}}{\norm{\vA \vu_{t}}_2},
\end{equation}
这类算法包括一个涉及与数据相关的矩阵 $\vA$
的“求幂”操作，以及一个强制执行所需结构类型的“投影”操作。\Cref{tab:power-method-summary}
总结了我们在本章中研究的遵循此结构的算法。读者可以体会到这种方法对不同类型的低维结构和不同损失（即PCA的二次损失和ICA的峰度型损失）的适用性，以及尽管关于这些算法的文献很广泛，但对于过完备字典学习却缺乏这样的算法。我们将为更多低维结构家族（特别是那些与深度学习盛行的应用相关的结构）开发幂法，视为本章提出的更重要（且开放）的研究问题之一。

我们在 \Cref{sec:mixture-and-dict}
中建立的几何子空间混合分布假设与更便于分析的稀疏字典假设之间的联系，在之前的工作中已被提及，特别是那些专注于广义主成分分析和子空间聚类等应用的工作，例如
\textcite{GPCA}
的工作。子空间混合假设将在整本书中继续发挥重要作用，既作为不同算法范式的分析测试用例，也作为推导不同深度网络架构的基础（如
\Cref{sec:LISTA} 中的 LISTA），但它可以扩展到更复杂的数据分布。

\section{练习与扩展}

\begin{exercise}\label{exercise:principal-components-derivation}
    证明，对于任何对称矩阵 \(\vA\)，问题 \(\max_{\vU \in \O(D,
    d)}\tr\left(\vU^{\top}\vA\vU\right)\) 的解是矩阵 \(\vU^{\star}\)，其列是
    \(\vA\) 的前 \(d\) 个单位特征向量。
\end{exercise}

\begin{exercise}\label{exercise:gaussian-rot-invar}
    设 $\vz \sim \cN(\Zero, \sigma^2 \vI)$ 是一个具有独立成分的高斯随机变量，每个成分的方差为 $\sigma^2$。
    证明对于任何正交矩阵 $\vQ$（即 $\vQ^\top \vQ =
    \vI$），随机变量 $\vQ \vz$ 的分布与 $\vz$ 相同。
    \textit{（提示：回忆高斯概率密度函数的公式，以及随机变量线性函数的密度公式。） }
\end{exercise}

\begin{exercise}\label{exercise:symmetry-identifiability}
    上面讨论的统计可识别性概念可以与模型类的\textit{对称性}（symmetries）联系起来，从而允许以纯确定性的方式理解估计，而无需任何统计假设。

    考虑模型 $\vX = \vU \vZ$，其中矩阵 $\vX, \vU, \vZ$
    具有兼容的尺寸。
    \begin{enumerate}
        \item 证明如果 $\vA$ 是任何兼容尺寸的方形可逆矩阵，那么在该模型下，对
            $(\vU \vA, \vA^{-1} \vZ)$ 也等于 $\vX$。我们称之为\textit{$\GL(d)$ 对称性}。
        \item 假设 $\vZ$ 的每一列都是来自公共统计模型 $\vz$
            的独立同分布观测值，此外该模型具有零均值和具有正方差的独立成分 $z_i$。
            证明对于任何方形可逆矩阵 $\vA$，如果 $\vA
            \vz$ 具有不相关的成分，那么 $\vA$ 可以写成 $\vD_1 \vQ \vD_2$，其中 $\vQ$
            是正交矩阵，$\vD_1$、$\vD_2$ 是对角矩阵。
            \textit{这把 ICA 中的“独立性”假设与“对称性破缺”效应联系起来，后者只允许尺度和旋转对称性。}
    \end{enumerate}

\end{exercise}

\begin{exercise}\label{exercise:whitening}
    考虑模型 $\vx = \vU \vz$，其中 $\vU \in \R^{D \times
    d}$（$D \geq d$）是固定的且秩为 $d$，$\vz$ 是零均值随机变量。设 $\vx_1, \dots, \vx_N$
    表示来自该模型的独立同分布观测值。
    \begin{enumerate}
        \item 证明矩阵 $\vX = [\vx_1, \dots, \vx_N]$ 的秩不大于
            $d$，因此存在一个标准正交矩阵 $\vV \in \R^{D \times d}$ 使得 $\vX
            = \vV \vY$，其中 $\vY \in \R^{d \times N}$。
            \textit{（提示：使用 PCA。）}
        \item 证明只要 $\Cov(\vz)$ 是非奇异的，\textit{白化矩阵} $(\vY
            \vY^\top)^{-1/2}\vY$
            在期望上就存在，并且它具有单位经验协方差。\footnote{特别地，可以在数学上证明，只要 $\vz$
            满足适当的集中不等式且 $N$ 足够大，这就足以保证白化矩阵以高概率存在。}
        \item 通过使用 $\vU$ 的奇异值分解，证明可以选择矩阵 $\vV$ 使得白化矩阵满足 $(\vY
            \vY^\top)^{-1/2} \vY =
            \vW [\vz_1, \dots, \vz_N]$，其中 $\vW$ 是标准正交矩阵。
    \end{enumerate}
\end{exercise}

\begin{exercise}\label{exercise:kurtosis-linearity-properties}
    设 $X$ 和 $Y$ 为零均值独立随机变量。
    \begin{enumerate}
        \item 证明 $\kurt(X + Y) = \kurt(X) + \kurt(Y)$。
        \item 对于任何 $\alpha \in \R$，证明 $\kurt(\alpha X) =
            \alpha^4 \kurt(X)$。
    \end{enumerate}
\end{exercise}

\begin{exercise}\label{exercise:sphere-calculus}
    设 $f : \R^d \to \R$ 为给定的二次连续可微目标函数。考虑球面约束优化问题
    \begin{equation}\label{eq:exercise-sphere-constrained-max}
        \max_{\norm{\vu}_2^2 = 1}\, f(\vu).
    \end{equation}
    在本练习中，我们将推导在 FastICA 推导中给出的表达式，用于通过梯度上升算法在球面上最大化峰度。
    这些表达式是流形上的微积分和优化丰富理论的特例，其中球面是一个特定的例子。对该领域进行深入的技术研究超出了我们的目的，因此我们只为感兴趣的读者提及两篇关键参考文献：Absil、Mahony
    和 Sepulchre 的开创性教科书
    \cite{Absil2009-nc}，以及 Boumal 最近的入门论著
    \cite{Boumal2023-rj}。
    \begin{enumerate}
        \item 对于任何作为 $\R^d$ 的可微子流形的约束集 $\cM$，在点 $\vu \in \cM$
            处的\textit{切空间}（tangent space），非正式地讲，是流形 $\cM$ 在点 $\vu$ 处的最佳局部线性近似。
            在一个重要的特例中，如果 $\cM$ 在 $\vu$ 处局部定义为函数 $F : \R^d \to \R$ 的水平集，即
            \begin{equation*}
                U \cap \cM = F^{-1}(\set{0})
            \end{equation*}
            对于某个包含 $\vu \in U$ 的开集 $U \subset \cM$，
            则 $\cM$ 在 $\vu$ 处的切空间可以通过微分计算：
            \begin{equation*}
                T_{\vu} \cM = \Ker(DF_{\vu}).
            \end{equation*}
            很容易看出，球面具有定义方程 $F(\vu) = \norm{\vu}_2^2 - 1$。
            利用这些事实证明，球面在 $\vu$ 处的切空间由下式给出
            \begin{equation*}
                T_{\vu} \bS^{d-1} = \set{\vv \in \R^d \given \ip{\vv}{\vu} = 0},
            \end{equation*}
            并且到该子空间的正交投影为
            $\vP_{\vu}^\perp = \vI - \vu\vu^\top$。
        \item 向量场
            \begin{equation}\label{eq:exercise-riemann-grad-sphere}
                \mathrm{grad}\, f(\vu) = \vP_{\vu}^\perp \nabla f
            \end{equation}
            被称为限制在球面上的函数 $f$ 的\textit{黎曼梯度}（Riemannian gradient）。
            优化问题
            \eqref{eq:exercise-sphere-constrained-max}
            的\textit{一阶最优性条件}可以用黎曼梯度表示：
            \begin{equation*}
                \mathrm{grad}\, f(\vu) = \mathbf{0}.
            \end{equation*}
            几何上，这表示 $f$ 在 $\vu$ 处的欧几里得梯度必须正交于球面在 $\vu$ 处的切空间。
            现在假设 $\vv \in \R^d$ 非零。使用一阶最优性条件证明
            \begin{equation*}
                \mathrm{proj}_{\bS^{d-1}}(\vv) \doteq
                \min_{\norm{\vu}_2^2 = 1}\, \norm{\vu - \vv}_2
                =
                \frac{\vv}{\norm{\vv}_2}，
            \end{equation*}
        \item 在 $\R^d$ 上的优化中，人们使用\textit{海森矩阵}（Hessian matrix）
            $\nabla^2 f(\vu)$ 来检查二阶最优性条件（以确定临界点是最大值点、最小值点还是鞍点）。
            通过像本练习第一部分那样对球面的黎曼梯度
            $\mathrm{grad}\, f(\vu)$ 关于 $\vu$
            求导，证明用于确定球面约束优化的二阶最优性条件的对应对象是\textit{黎曼海森矩阵}（Riemannian Hessian），定义为
            \begin{equation}\label{eq:exercise-riemann-hess-sphere}
                \mathrm{Hess}\, f(\vu) = \vP_{\vu}^\perp \left(
                    \nabla^2 f(\vu) - \ip{\nabla f(\vu)}{\vu} \vI
                \right) \vP_{\vu}^\perp.
            \end{equation}
    \end{enumerate}
\end{exercise}

\begin{exercise}\label{exercise:kurtosis-sphere-landscape}
    在本练习中，我们简述了文献中称为\textit{景观分析}（landscape analysis）的论证，用于球面约束的总体峰度最大化问题
    \eqref{eq:kurtosis-maximization-sphere-population-simple}。我们将证明，当至少存在一个具有正峰度的独立成分时，其全局最大值点确实会导致恢复字典
    $\vU$ 的一列。
    为简单起见，我们将假设对于每个 $i = 1, \dots, d$，$\kurt(z_i) \neq 0$。
    \begin{enumerate}
        \item 使用 \Cref{exercise:sphere-calculus} 第1部分的结果，证明
            \eqref{eq:kurtosis-maximization-sphere-population-simple} 的一阶最优性条件为
            \begin{equation}\label{eq:kurtosis-sphere-landscape-stationarity}
                \left(\sum_{i=1}^d \kurt(z_i) w_i^4 \right)
                \vw = \kurt(\vz) \hada \vw^{\hada 3},
            \end{equation}
            其中峰度是逐元素计算的，$\hada$
            表示向量的逐元素乘法，$\vw^{\hada 3}$ 表示其参数的逐元素立方。
        \item 证明具有单位范数且满足 \eqref{eq:kurtosis-sphere-landscape-stationarity}
            的向量 $\vw$ 均采用以下形式。
            设 $S^+ = \set{i \in [d] \given \kurt(z_i) > 0}$，且
            $S^- = \set{i \in [d] \given \kurt(z_i) < 0}$。
            设 $S$ 为 $S^+$ 或 $S^-$ 的子集。
            那么
            \begin{equation}\label{eq:kurtosis-sphere-landscape-cps}
                \vw_S = \sum_{i \in S} \pm \sqrt{\frac{1}{\kurt(z_i)
                \sum_{j \in S} \frac{1}{\kurt(z_j)}}} \ve_i
            \end{equation}
            满足 \eqref{eq:kurtosis-sphere-landscape-stationarity}，
            其中 $\ve_i$ 是在第 $i$ 个位置为 $1$ 且在其他位置为 $0$ 的向量，$\pm$ 符号表示选择正号或负号。
        \item 假设至少存在一个 $i$ 使得
            $\kurt(z_i) > 0$。使用 \Cref{exercise:sphere-calculus} 第2部分的结果，证明
            \eqref{eq:kurtosis-maximization-sphere-population-simple}
            目标的唯一局部最大值点是带符号的单稀疏向量 $\pm \ve_i$（其中 $i \in
            S^+$）。得出结论，\eqref{eq:kurtosis-maximization-sphere-population-simple}
            的全局最大值点是对应于具有最大峰度成分的带符号单稀疏向量。 %
            \textit{（提示：计算每个临界点处黎曼海森矩阵
            \eqref{eq:exercise-riemann-hess-sphere} 的正负特征值的数量。）}
        \item 现在假设对于每个 $j =1,
            \dots, d$，$\kurt(z_j) <
            0$。这对应于峰度最大化问题的“过度收缩”（over-deflated）实例。再次使用
            \Cref{exercise:sphere-calculus} 第2部分的结果，证明
            \eqref{eq:kurtosis-maximization-sphere-population-simple}
            目标的唯一局部最大值点是带符号的稠密向量 $\sum_{i=1}^d \pm \ve_i$。
            这表明优化公式
            \eqref{eq:kurtosis-maximization-sphere-population-simple}
            不能被简单地套用。
    \end{enumerate}
\end{exercise}

\begin{exercise}\label{exercise:orthogonal-group-calculus}
    本练习遵循 \Cref{exercise:sphere-calculus} 中介绍的结构和形式，但将其应用于正交群 $\O(d)
    = \set{\vU \in \R^{d \times d} \given
    \vU^\top \vU = \vI}$。
    有关必要的概念背景，请参阅 \Cref{exercise:sphere-calculus} 的描述；只要记住矩阵上的相关内积是
    $\ip{\vX}{\vY} = \tr(\vX^\top \vY)$，该形式就同样适用于环境空间为 $d
    \times d$ 矩阵集合的情况。
    关于正交群优化的一个极好的通用参考文献是 Edelman、Arias 和 Smith 的著作 \cite{Edelman1998-lg}。

    设 $f : \R^{d\times d} \to \R$ 为给定的二次连续可微目标函数。考虑正交约束优化问题
    \begin{equation}\label{eq:exercise-orthogonal-group-constrained-max}
        \max_{\vQ^\top \vQ = \vI}\, f(\vQ).
    \end{equation}
    \begin{enumerate}
        \item  很容易看出，正交群具有定义方程 $F(\vQ) = \vQ^\top \vQ = \vI$。
            利用这一事实证明，正交群在 $\vQ$ 处的切空间由下式给出
            \begin{equation*}
                T_{\vQ} \O(d) = \set{\vQ \vOmega \in \R^{d\times d}
                \given \vOmega^\top = - \vOmega},
            \end{equation*}
            并且到该子空间的正交投影为
            \begin{equation*}
                \cP_{T_{\vQ}\O(d)}(\vDelta) =  \vQ \Skew(\vQ^\top \vDelta),
            \end{equation*}
            其中 $\Skew(\vDelta) = \tfrac{1}{2} ( \vDelta -
            \vDelta^\top)$ 是到斜对称矩阵集合的正交投影。
            向量场
            \begin{equation}\label{eq:exercise-riemann-grad-orthogonal-group}
                \mathrm{grad}\, f(\vQ) = \cP_{T_{\vQ} \O(d)} \left(
                \nabla f(\vQ) \right)
            \end{equation}
            被称为限制在正交群上的函数 $f$ 的\textit{黎曼梯度}。
            优化问题
            \eqref{eq:exercise-orthogonal-group-constrained-max}
            的\textit{一阶最优性条件}可以用黎曼梯度表示：
            \begin{equation*}
                \mathrm{grad}\, f(\vQ) = \mathbf{0}.
            \end{equation*}
        \item 通过像本练习第一部分那样对正交群的黎曼梯度
            $\mathrm{grad}\, f(\vQ)$ 关于 $\vQ$ 求导，证明\textit{黎曼海森矩阵}由下式给出
            \begin{equation}\label{eq:exercise-riemann-hess-orthogonal-group}
                \mathrm{Hess}\, f(\vQ) = \cP_{T_{\vQ}\O(d)} \left(
                    \nabla^2 f(\vQ) - \Symm(\vQ^\top \nabla f(\vQ)) \kron \vI
                \right) \cP_{T_{\vQ}\O(d)},
            \end{equation}
            其中 $\Symm(\vDelta) = \tfrac{1}{2}(\vDelta +
            \vDelta^\top)$ 表示到对称矩阵集合的正交投影，$\kron$ 表示矩阵的克罗内克积（Kronecker product）。
            请注意，将前一个表达式中出现的算子解释为\textit{在 ${d \times d}$
            矩阵上的线性变换}，而\textbf{不是} $d \times d$ 矩阵本身。
            优化问题
            \eqref{eq:exercise-orthogonal-group-constrained-max}
            的\textit{二阶最优性条件}可以用黎曼海森矩阵表示：
            \begin{equation*}
                \mathrm{Hess}\, f(\vQ) \preceq \mathbf{0}.
            \end{equation*}
            对于最小化问题，符号反转。

            \textit{（提示：关键是操作计算以获得尽可能紧凑的形式
                \eqref{eq:exercise-riemann-hess-orthogonal-group}。为此，利用克罗内克积的以下同构：如果
                $\vA$、$\vX$ 和 $\vB$ 是兼容尺寸的矩阵，那么有
                \begin{equation*}
                    (\vB^\top \kron \vA) \Vec(\vX) = \Vec(\vA \vX \vB),
                \end{equation*}
                其中 $\Vec(\vX)$ 表示将矩阵参数的列“从左到右”堆叠成一个向量。我们在
                \eqref{eq:exercise-riemann-hess-orthogonal-group}
            中使用这种同构，以便以规范的方式将两个矩阵的克罗内克积定义为矩阵上的算子。）}

        \item 现在假设 $\vX \in \R^{d\times d}$ 是满秩的。在本练习的这一部分和下一部分中，我们考虑
            $\vX$ 到正交群的投影：
            \begin{equation}\label{eq:projection-onto-ogrp-defn}
                \mathrm{proj}_{\O(d)}(\vX) \doteq
                \min_{\vQ \in \O(d)
                }\, \norm{\vQ - \vX}_F^2.
            \end{equation}
            我们将证明该问题的解由下式给出
            \begin{equation*}
                \mathrm{proj}_{\O(d)}(\vX)
                =
                \vU \vV^\top,
            \end{equation*}
            其中 $\vX = \vU \vS \vV^\top$ 是 $\vX$ 的奇异值分解。

            \begin{enumerate}
                \item 使用一阶和二阶最优性条件，证明
                    \eqref{eq:projection-onto-ogrp-defn} 的每个局部最小值点 $\vQ$ 都满足
                    \begin{align*}
                        \left( \vQ^\top \vX \right)^\top &= \vQ^\top \vX, \\
                        \vQ^\top \vX &\succeq \mathbf{0}.
                    \end{align*}
                    \textit{（提示：当克罗内克积的两个参数中的一个固定时，利用其在另一个参数上的线性。）}
                \item 使用这些条件，论证在 \eqref{eq:projection-onto-ogrp-defn}
                    的每个局部最小值点 $\vQ$ 处，都有 $\vQ^\top \vX = (\vX^\top \vX)^{1/2}$。
                    \textit{（提示：使用 %
                        线性代数中的事实，即如果 $\vS \succeq
                    \mathbf{0}$ 是对称半正定矩阵，那么 $(\vS^\top\vS)^{1/2} = \vS$。）}
                \item 使用奇异值分解 $\vX =
                    \vU \vS \vV^\top$，得出结论
                    \begin{equation*}
                        \vU \vV^\top
                        =
                        \proj_{\O(d)}(\vX).
                    \end{equation*}
            \end{enumerate}
    \end{enumerate}
\end{exercise}

\end{document}

